\documentclass[12pt]{report}
\usepackage[utf8]{inputenc}
\usepackage[spanish]{babel}
\usepackage{amsmath, amssymb}

\usepackage{array}

\usepackage{graphicx}
\usepackage{lmodern}
\usepackage{fancyhdr}
\usepackage{geometry}
\usepackage{titlesec}
\usepackage{parskip}
\usepackage{ragged2e}
\usepackage{hyperref} 
\hypersetup{colorlinks=true, urlcolor=blue, linkcolor=black, citecolor=blue}
\newcommand{\citeref}[2]{\hyperlink{ref:#1}{#2}}
\usepackage{tabularx}
\usepackage{booktabs}
\usepackage{caption}
\usepackage{float}
\usepackage{enumitem}
\usepackage{caption}  
\usepackage{longtable} 
\captionsetup[table]{name=Tabla}
\geometry{margin=3cm}

% Configuración de espaciado en listas
\setlist[itemize]{itemsep=0.3cm, parsep=0.3cm}
\setlist[enumerate]{itemsep=0.3cm, parsep=0.3cm}

% Mejora de tipografía de subsecciones
\titleformat{\section}{\normalfont\Large\bfseries}{\thesection}{0.5em}{}
\titleformat{\subsection}{\normalfont\large\bfseries}{\thesubsection}{0.5em}{}
\titleformat{\subsubsection}{\normalfont\bfseries}{\thesubsubsection}{0.5em}{}
\setlength{\headheight}{14.5pt}
\pagestyle{fancy}
\fancyhf{}
\cfoot{\thepage}
\renewcommand{\headrulewidth}{0pt}
\renewcommand{\footrulewidth}{0pt}
\renewcommand{\baselinestretch}{1.25}
\setlength{\parindent}{1.25cm}
\setlength{\parskip}{0.1cm}

\titleformat{\chapter}{\normalfont\Large\bfseries}{\thechapter.\quad}{0pt}{}
\titlespacing*{\chapter}{0pt}{0pt}{1em}

% Números arábigos para capítulos y figuras, arábigos para secciones y subsecciones
\renewcommand{\thechapter}{\arabic{chapter}}
\renewcommand{\thesection}{\thechapter.\arabic{section}}
\renewcommand{\thesubsection}{\thechapter.\arabic{section}.\arabic{subsection}}
\renewcommand{\thesubsubsection}{\thechapter.\arabic{section}.\arabic{subsection}.\arabic{subsubsection}}
\renewcommand{\thefigure}{\arabic{chapter}.\arabic{figure}}
\usepackage{caption}

\captionsetup[figure]{name=Figura}
% Cambiar etiquetas de tablas de "Cuadro" a "Tabla"
\addto\captionsspanish{% spanish language
  \renewcommand{\tablename}{Tabla}%
  \renewcommand{\listtablename}{Índice de tablas}%
}

\begin{document}
\justifying

\begin{titlepage}
\centering

% -------- LOGOS --------
\begin{minipage}{0.45\textwidth}
    \flushleft
    \includegraphics[height=2.3cm]{fig/LogoEscuela.png}
\end{minipage}
\begin{minipage}{0.45\textwidth}
    \flushright
    \includegraphics[height=2.3cm]{fig/LogoUCN_acentuado.png}
\end{minipage}

\vspace{1.3cm}

% -------- TITULO (UN SOLO BLOQUE) --------
{\LARGE \bfseries 
Modelo de Soporte a la Toma de Decisiones en la Agricultura Basado en Simulación Multiparadigma
\par}

\vspace{0.4cm}

% -------- SUBTITULO --------
{\normalsize \itshape Informe Final Capstone de Investigación \par}

\vspace{0.7cm}

% -------- UNIVERSIDAD --------
{\small Universidad Católica del Norte \par}
{\small Escuela de Ingeniería \par}
{\small Sede Coquimbo \par}

\vspace{0.8cm}

% -------- INTEGRANTES --------
{\normalsize \textbf{Integrantes} \par}
\vspace{0.2cm}
{\small
Claudio Malebrán Cabezas \\
Daniel Morales González \\
Antonio Pereira Vergara
}

\vspace{1cm}

% -------- Tutor --------
{\normalsize \textbf{Tutores} \par}
\vspace{0.2cm}
{\small 
Dr. Carlos Monardes Concha \\
Ing. Agustin Olivares Soto}



\vspace{0.5cm} 

{\normalsize \textbf{Corrector} \par}
\vspace{0.2cm}
{\small Dr. Paulo Alves Junior\par}

\vspace{1cm} 
% -------- FECHA --------
{\small Coquimbo, Chile \par}
{\small \today \par}

\end{titlepage}
\tableofcontents

\listoffigures

\listoftables

% --- Resumen en español ---
\newpage
\selectlanguage{spanish}
\section*{Resumen}

En el presente, tras una creciente crisis hídrica que afecta a diversas zonas agrícolas del norte del país exige nuevas herramientas que permitan anticipar escenarios de incertidumbre y optimizar el uso de recursos. La caída sostenida de las aguas superficiales y el uso de acuíferos han puesto en peligro la continuidad productiva, haciendo necesario modelos que utilicen dinámicas ambientales y decisiones operativas. Dado este contexto, se desarrolla un modelo de simulación multiparadigma orientado a generar evidencia empírica para ayudar en la toma de decisiones agrícolas, a través de la combinación de dinámica de sistemas para evaluar el comportamiento de la oferta hídrica y eventos discretos para representar el sector operativo en la gestión de cultivos. Este modelo se aplica en el sector agrícola de Pan de Azúcar, en la Región de Coquimbo, caracterizado por una crisis hídrica estructural marcada por la severa reducción de aguas superficiales y el descenso crítico del acuífero El Culebrón.

%La implementación práctica del modelo se materializa en un Sistema de Soporte a la Decisión (DSS) a escala predial. Al contrastar los resultados obtenidos con los métodos tradicionales basados en experiencia empírica, la investigación evidencia que la anticipación de escenarios permite un uso más eficiente del recurso hídrico, contribuyendo a mitigar pérdidas y fortalecer la resiliencia productiva de los agricultores.

%El presente trabajo de investigación aborda la crisis hídrica estructural en el sector agrícola de Pan de Azúcar, Región de Coquimbo, escenario caracterizado por la severa restricción de aguas superficiales y el descenso crítico en los niveles del acuífero local El Culebron.El objetivo principal de este estudio es desarrollar un modelo de simulación multiparadigma que genere evidencia empírica para la toma de decisiones agrícolas frente a la incertidumbre.Para ello, el modelo integra la dinámica de sistemas que evalúa el comportamiento de la oferta de agua y eventos discretos para planificar las restricciones operativas de los cultivos. La materialización práctica de este modelo se consolida a través de un Sistema de Soporte a la Decisión (DSS) a escala predial. Al contrastar la información arrojada por este modelo con los métodos tradicionales de gestión basados en la experiencia empírica, la investigación busca generar evidencia que demuestre que la anticipación de escenarios conlleva a un mejor uso del recurso hídrico, permitiendo mitigar pérdidas y dotar de resiliencia productiva a los agricultores.
%---\textbf{Palabras clave:} agricultura de precisión, desarrollo sostenible, sensores.

\vspace{0.5cm}

% --- Abstract en inglés ---
\selectlanguage{english}
\section*{Abstract}
    \textit{This research study addresses the structural water crisis in the agricultural sector of Pan de Azúcar, Coquimbo Region, a setting characterized by severe surface water shortages and a critical decline in the levels of the local El Culebron aquifer.} \textit{The main objective of this study is to develop a multi-paradigm simulation model that generates empirical evidence for agricultural decision-making in the face of uncertainty. To this end, the model integrates system dynamics, which evaluates the behavior of water supply, and discrete events to plan for operational crop constraints.} \textit{The practical implementation of this model is consolidated through a Decision Support System (DSS) at the farm level. By comparing the information generated by this model with traditional management methods based on empirical experience, the research seeks to generate evidence demonstrating that anticipating scenarios leads to better use of water resources, thereby mitigating losses and providing farmers with productive resilience.}

%---\textbf{Keywords:} precision agriculture, sustainable development, sensors.

\chapter{Contextualización del fenómeno de investigación}
\section{Introducción}

La gestión del agua constituye uno de los principales desafíos para la sostenibilidad agrícola y la seguridad alimentaria a nivel mundial. De acuerdo con el Informe Mundial de las Naciones Unidas sobre el Desarrollo de los Recursos Hídricos elaborado por \citeref{un2023}{UNESCO (2024)}, la agricultura representa aproximadamente el 70\% del consumo de agua dulce a escala global, evidenciando la elevada presión que ejerce este sector sobre los recursos hídricos disponibles.

Esta situación también es respaldada por la literatura científica. Mekonnen y Hoekstra (2016) identifican a la agricultura como el principal consumidor de agua azul a nivel mundial, mientras que Palma y Tomaz (2024) señalan que el incremento de la demanda alimentaria, la expansión de la superficie cultivada y los efectos del cambio climático continuarán intensificando la competencia por los recursos hídricos durante las próximas décadas.

En Chile, esta problemática adquiere una relevancia aún mayor debido a la prolongada megasequía que afecta al país desde el año 2009 y a la elevada variabilidad hidrológica característica de gran parte del territorio nacional. \citeref{rivera2025}{Rivera y Donoso (2025)} señalan que la gestión de la sequía en Chile ha estado históricamente orientada hacia una respuesta reactiva frente a las crisis hídricas, enfatizando la necesidad de avanzar hacia estrategias preventivas sustentadas en sistemas de monitoreo, planificación y apoyo a la toma de decisiones. Los autores destacan además que la complejidad del fenómeno exige integrar dimensiones hidrológicas, ambientales, económicas, institucionales y sociales para lograr una gestión sostenible del recurso hídrico.

Desde el punto de vista normativo, el Código de Aguas establece que el agua constituye un bien nacional de uso público cuyo aprovechamiento debe compatibilizar el desarrollo productivo con la preservación ecosistémica y la satisfacción del consumo humano, otorgando prioridad a este último en situaciones de escasez hídrica (\citeref{codigo1981}{Código de Aguas}, DFL N° 1.122; Ley N° 21.435, 2022). En consecuencia, la disponibilidad efectiva del recurso para la agricultura no depende únicamente de las condiciones hidrológicas, sino también de las decisiones adoptadas por las organizaciones de usuarios del agua y del marco regulatorio vigente.

\begin{figure}[h]
    \centering
    \includegraphics[width=0.8\textwidth]{fig/caudal rio.png}
    \caption{Déficit hídrico por cuenca en Chile}
    \caption*{\footnotesize Fuente: Boletín Hidrométrico \citeref{dga2026}{DGA (2026)}}
    \label{fig:deficit_hidrico_chile}
\end{figure}

Una de las regiones donde esta problemática se manifiesta con mayor intensidad corresponde a la Región de Coquimbo, la cual permanece bajo decreto de escasez hídrica desde el año 2008. Esta condición ha generado importantes repercusiones sobre la actividad agrícola y la disponibilidad de agua para riego. En particular, como se aprecia en la Figura~\ref{fig:deficit_hidrico_chile}, la cuenca del río Elqui presenta un déficit hídrico cercano al 82\% respecto de sus promedios históricos (Boletín Hidrométrico N°574, \citeref{dga2026}{DGA, 2026}), situación consistente con los efectos de la megasequía descritos por \citeref{rivera2025}{Rivera y Donoso (2025)} para la zona centro-norte del país.

Dentro de las áreas más afectadas se encuentra el valle de Pan de Azúcar, donde la mayor parte de los agricultores depende del canal Bellavista como principal fuente de abastecimiento hídrico. De acuerdo con lo señalado por Manuel Domínguez, ingeniero civil asesor de la Junta de Vigilancia del Canal Bellavista, durante una reunión sostenida con regantes del sector, ``el canal posee una capacidad de riego aproximada de 13.000 hectáreas; sin embargo, debido a las condiciones actuales de sequía, solo se encuentran en producción cerca de 4.800 hectáreas'' (Domínguez, comunicación personal, 26 de marzo de 2026), evidenciando la importante reducción experimentada por la superficie agrícola efectiva.

\begin{figure}[h]
    \centering
    \includegraphics[width=1\textwidth]{fig/mapa.png}
    \caption{Sistema de abastecimiento hídrico en el Valle de Pan de Azúcar}
    \caption*{\footnotesize Fuente: Elaboración propia mediante Google Earth}
    \label{fig:mapa_abastecimiento}
\end{figure}

Como se observa en la Figura~\ref{fig:mapa_abastecimiento}, el área agrícola del valle de Pan de Azúcar depende directamente del río Elqui mediante el canal Bellavista, evidenciando la estrecha relación existente entre la disponibilidad hídrica de la cuenca y la producción agrícola del territorio. Esta dependencia incrementa la vulnerabilidad del sistema productivo frente a períodos prolongados de escasez hídrica y frente a las decisiones de administración del recurso adoptadas por las organizaciones de usuarios del agua.

En este contexto, la actividad agrícola se encuentra además condicionada por la priorización legal del recurso para consumo humano y saneamiento, lo que implica la aplicación de medidas restrictivas sobre el uso agrícola del agua. 

Entre estas medidas se encuentra el desmarque de las acciones de agua, mecanismo mediante el cual se ajusta el porcentaje efectivo de derechos de aprovechamiento que pueden ser utilizados durante cada temporada de riego, en función de la disponibilidad hídrica proyectada por la Junta de Vigilancia. En escenarios críticos, cuando el volumen almacenado en el embalse Puclaro desciende por debajo de aproximadamente 8 millones de metros cúbicos, incluso puede llegar a suspenderse completamente la entrega de agua para riego. 

Esta situación estuvo próxima a materializarse durante el año 2024, cuando el embalse registró un volumen cercano a 8,7 millones de metros cúbicos, lo que motivó la aplicación de un desmarque del 10\%.

Bajo este escenario de elevada incertidumbre hidrológica, disponibilidad variable del recurso y restricciones operacionales para la agricultura, resulta fundamental desarrollar herramientas que permitan apoyar la toma de decisiones mediante la evaluación anticipada de distintos escenarios de disponibilidad hídrica. 

En este contexto, la simulación constituye una alternativa particularmente adecuada para integrar la incertidumbre asociada a la oferta de agua, el comportamiento del sistema productivo y las consecuencias económicas de distintas estrategias de cultivo, contribuyendo así a fortalecer la resiliencia hídrica y la planificación agrícola en sistemas hortícolas de la Región de Coquimbo.

\section{Pregunta de investigación}


¿Puede la simulación generar evidencia que reduzca la incertidumbre en la toma de decisiones en el sector agrícola de pan de azúcar, en contraste con prácticas comunes en escenarios de escasez hídrica?


\section{Objetivo general}
Generar evidencia empírica replicable sobre la gestión del agua que permita mejorar la toma de decisiones  en el sector productivo agrícola utilizando herramientas de las ciencias de la ingeniería.

\section{Objetivos específicos}
\begin{enumerate}
    \item Caracterizar la gestión actual del recurso hídrico en sistemas agrícolas, identificando prácticas de riego y principales fuentes de agua. 
    \item Analizar los procesos de toma de decisiones de los agricultores, identificando las variables consideradas y las principales limitaciones.
    \item Desarrollar un modelo  que permita generar evidencia para la toma de decisiones en un sistema productivo agrícola.
    \item Materializar de manera práctica el modelo con el desarrollo de un  DSS que pueda generar impacto en la comunidad agrícola.
    \item Evaluar el impacto del modelo y del DSS en la eficiencia del uso del agua y en la generación de valor económico en el sector agrícola.

\end{enumerate}  

\section{Marco conceptual}

\vspace{0.3cm}

\paragraph{Evapotranspiración de referencia (ETo):} 
Corresponde a la demanda evaporativa de la atmósfera sobre una superficie de referencia (pasto), bajo condiciones estándar. Se utiliza como base para estimar el consumo hídrico de distintos cultivos y es fundamental en la planificación del riego.

\vspace{0.2cm}

\paragraph{Evapotranspiración del cultivo (ETc):} 
Representa la cantidad de agua que un cultivo específico requiere, considerando sus características y estado de desarrollo. Se obtiene a partir de la ETo ajustada mediante el coeficiente de cultivo (Kc), permitiendo estimar necesidades reales de riego.

\vspace{0.2cm}

\paragraph{Coeficiente de cultivo (Kc):} 
Es un factorde ajuste la evapotranspiración de referencia (ETo) que la escala a la evapotranspiración del cultivo (ETc). Varía según el tipo de cultivo y su etapa fenológica, siendo clave para determinar la demanda hídrica precisa.
\begin{equation}
ET_c = ET_0 \cdot K_c
\end{equation}

\vspace{0.2cm}

\paragraph{Sistema de riego:} 
Corresponde al conjunto de técnicas y tecnologías utilizadas para aplicar agua a los cultivos de manera controlada. Incluye métodos como riego por gravedad, aspersión y goteo, los cuales difieren en su eficiencia, uniformidad de aplicación y pérdidas por evaporación o escorrentía. La tecnificación del sistema de riego influye directamente en la eficiencia del uso del agua y en la sostenibilidad del sistema agrícola.

\vspace{0.2cm}

\paragraph{Umbrales de riego:} 
Corresponden a niveles de humedad del suelo que delimitan estados relevantes para la disponibilidad de agua para las plantas y la gestión del riego. Estos umbrales permiten interpretar cuándo el suelo se encuentra en condiciones óptimas, de estrés o de déficit hídrico. Su valor depende directamente de las características físicas del suelo, tales como textura, estructura y capacidad de retención de agua.

\vspace{0.2cm}

\paragraph{Capacidad de campo (CC):} 
Corresponde a un umbral superior de humedad del suelo, asociado al nivel máximo de agua que el suelo puede retener después del drenaje por gravedad. En este punto, el agua se encuentra disponible para las plantas sin generar saturación, constituyendo una condición óptima para el desarrollo del cultivo.

\vspace{0.2cm}

\paragraph{Punto de marchitez permanente (PMP):} 
Corresponde a un umbral inferior de humedad del suelo, en el cual las plantas ya no pueden extraer agua, incluso bajo condiciones ambientales favorables. Este nivel define el límite crítico de disponibilidad hídrica, asociado a condiciones de estrés severo e irreversible.


\vspace{0.2cm}

\paragraph{Desmarques:} 
Mecanismo de distribución proporcional del agua en escenarios de escasez, donde cada usuario recibe un porcentaje de su derecho de agua según la disponibilidad total; por ejemplo, con un desmarque del 15\%, se entrega solo el 15\% del volumen asignado originalmente.

\vspace{0.2cm}

\paragraph{Pérdidas por conducción:} 
Corresponden a las pérdidas físicas de volumen de agua asociadas al transporte del recurso a través de la infraestructura hidráulica, desde un punto A de captación o bocatoma hasta un punto B de entrega predial, calculadas como una función de la distancia recorrida y del tipo de revestimiento del canal.

\vspace{0.2cm}

\paragraph{Capacidad de agua almacenada:} 
Corresponde al volumen máximo de agua que puede ser retenido en un sistema, como el suelo, un embalse o un acuífero, y que está disponible para su uso posterior en riego u otras actividades.

\vspace{0.2cm}

\paragraph{CEAZA:} 
Institución que monitorea y entrega información climática y ambiental en zonas áridas, especialmente en la región de Coquimbo, facilitando datos relevantes para la gestión hídrica y agrícola.

\vspace{0.2cm}

\paragraph{Dirección General de Aguas (DGA):} 
Organismo estatal encargado de la gestión, regulación y monitoreo de los recursos hídricos en Chile, incluyendo información sobre caudales, derechos de agua y disponibilidad hídrica.

\vspace{0.2cm}

\paragraph{Acuífero:} 
Formación geológica que almacena agua subterránea, la cual puede ser extraída mediante pozos y utilizada para riego u otros fines productivos.

\vspace{0.2cm}

\paragraph{Acciones de agua (derechos de agua):} 
Son derechos legales que permiten el uso de un determinado volumen o caudal de agua, ya sea de fuentes superficiales o subterráneas, y que regulan el acceso al recurso hídrico.

\vspace{0.2cm}

\paragraph{Canal de riego:} 
Infraestructura hidráulica utilizada para transportar agua desde fuentes naturales o artificiales hacia las zonas de cultivo.

\vspace{0.2cm}

\paragraph{Turnos de agua:} 
Sistema de distribución del agua en el cual los usuarios acceden al recurso en horarios o periodos definidos, permitiendo gestionar la escasez y asegurar una asignación equitativa.
\chapter{Revisión de la literatura}

\section{Marco teórico}

\small 
\begin{longtable}{p{3cm} p{3cm} p{5cm} p{4cm}}

% --- Encabezado de la primera página ---
\caption{Tabla resumen del Marco Teórico} \\
\hline
\textbf{Referencia / Documento} & \textbf{Enfoque Principal} & \textbf{Conceptos Clave} & \textbf{Relevancia para la Investigación}\\ \hline
\endfirsthead

% --- Encabezado de las siguientes páginas ---
\caption[]{Tabla de resumen(continuación)} \\
\hline
\textbf{Referencia / Documento} & \textbf{Enfoque Principal} & \textbf{Conceptos Clave} & \textbf{Relevancia para la Investigación}\\ \hline
\endhead

% --- Pie de página de la tabla (excepto la última) ---
\hline
\multicolumn{4}{r}{{Continúa en la siguiente página...}} \\
\endfoot

% --- Pie de página final ---
\hline
\endlastfoot

% --- Datos de la tabla ---
FAO Irrigation and Drainage Paper No. 56 & Cálculo técnico de la demanda hídrica de los cultivos & Evapotranspiración del cultivo ((ETc)), coeficiente de cultivo ((Kc)) y metodología de coeficientes duales & Proporciona la base matemática para estimar las necesidades reales de riego según la etapa fenológica y el clima\\ \hline

Droughts in Chile (Rivera, Donoso, 2025) & Análisis de la sequía desde dimensiones espaciales, temporales y económicas & Comportamiento de agentes económicos, brechas de información y efectos económicos directos de la escasez & Ayuda a comprender cómo la falta de datos limita la toma de decisiones eficientes y cómo responden los agricultores ante la crisis\\ \hline

DFL 1122 Código de Aguas & Marco legal y regulatorio de la gestión del agua en Chile & Aguas superficiales y subterráneas, prioridad para consumo humano, "desmarques" y zonas de restricción & Define las restricciones operativas y legales (como la reducción proporcional de derechos) que enfrentan los agricultores en escenarios de escasez\\ \hline

Dingre \& Gorantiwar (2020) & Determinación de requerimientos hídricos locales mediante balance de agua en suelo & Coeficiente de cultivo ((Kc)) específico, etapas de crecimiento y evapotranspiración real & Demuestra que los valores globales de FAO-56 pueden sobreestimar la demanda local hasta en un 16.6\%, subrayando la necesidad de calibración regional\\ \hline

Jiménez-Morales et al. (2014) & Modelaje matemático de la respuesta del crecimiento a la fertilización & Funciones de producción, biomasa (PSPA), área foliar e interacción de factores & Provee modelos de simulación para predecir el crecimiento del cultivo ante variables de manejo orgánico y mineral\\ \hline

Cai , Rosegrant (2004) & Selección de tecnología de riego bajo incertidumbre hidrológica & Programación estocástica de dos etapas, uniformidad de distribución (CUC) y análisis de aversión al riesgo & Ofrece un marco para equilibrar el riesgo de pérdida de beneficios en años secos frente al exceso de costo de capital en años húmedos\\\hline

Kuschel-Otárola et al. (2020) & Simulación de la Eficiencia del Uso del Agua (WUE) con diferentes estrategias & Modelo AquaCrop-OS, programación óptima (GET-OPTIS) y parámetros hidráulicos del suelo & Valida que la programación óptima maximiza la rentabilidad y la WUE al considerar la interacción suelo-cultivo-clima\\\hline




\end{longtable}




\newpage
\section{Estado del arte}


\small 
\begin{longtable}{p{2.8cm} p{2.8cm} c p{4cm} p{3cm}}

% --- Encabezado de la primera página ---
\caption{Tabla resumen de las investigaciones recientes y sus aportes} \\
\hline
\textbf{Investigación} & \textbf{Autor} & \textbf{Año} & \textbf{Objetivo} & \textbf{Revista científica } \\ \hline
\endfirsthead

% --- Encabezado de las siguientes páginas ---
\caption[]{Tabla de resumen(continuación} \\
\hline
\textbf{Investigación} & \textbf{Autor} & \textbf{Año} & \textbf{Objetivo} & \textbf{Revista científica} \\ \hline
\endhead

% --- Pie de página de la tabla (excepto la última) ---
\hline
\multicolumn{5}{r}{{Continúa en la siguiente página...}} \\
\endfoot

% --- Pie de página final ---
\hline
\endlastfoot

% --- Datos de la tabla ---
An agent-based model of farmer decision making & El Fartassi et al. & 2025 & Simular las decisiones de riego de los agricultores y su adaptación ante la escasez y salinidad del agua en regiones áridas & Agricultural Water Management \\ \hline

Effective Management of Water Resources Problems in Irrigated Agriculture & Ajay Singh & 2024 & Evaluar estrategias para mitigar problemas de salinización y anegamiento en agricultura de regadío mediante modelos de simulación & Water Resources Management (Springer) \\ \hline

Evaluación de los procesos de siembra y corte de flores & Osorio Gutierrez et al. & 2021 & Optimizar la rentabilidad de una empresa floricultora mediante simulación de eventos discretos para reducir retrasos & Prospectiva \\ \hline

Simulation-optimization based real-time irrigation scheduling & JLi et al. & 2023 & Desarrollar un método interactivo hombre-máquina para la programación de riego en tiempo real integrando asimilación de datos & Agricultural Water Management \\ \hline

Assessment of meteorological and agricultural droughts & Kaur et al. & 2022 & Monitorear sequías y su impacto en los niveles de agua subterránea mediante sensores remotos en la India & Agricultural Water Management \\ \hline

A Revisit of Internet of Things Technologies in Smart Agriculture & Rehman et al. & 2022 & Revisar las tecnologías IoT y estrategias de control para optimizar el rendimiento de los cultivos y el ahorro de recursos
 & Agronomy \\ \hline

Towards a Decision-Making Approach of Sustainable Water Management & Alitane et al. & 2022 & Cuantificar el balance hídrico en la cuenca del R’Dom (Marruecos) usando SWAT+ para apoyar la gestión sostenible & Sustainability \\ \hline

TWater balance model (WBM) v.1.0.0 & Grogan et al. & 2022 & Presentar un modelo hidrológico global de código abierto capaz de rastrear el origen y el uso del agua & Geoscientific Model Development \\ \hline

Integration of smart sensors and IOT in precision agriculture& Mansoor et al. & 2022 & Examinar las tendencias y desafíos de los sensores inteligentes e IoT para mejorar la eficiencia en la agricultura de precisión & Frontiers in Plant Science \\ \hline

Reinforce lake water balance component estimations & Mahmoodi et al. & 2025 & Integrar composiciones de isótopos de agua con modelos hidrológicos para mejorar la precisión del balance hídrico en lagos & Hydrology and Earth System Sciences \\ \hline

Monitoring the physical processes driving mass loss of Tapado Glacier & Ayala et al. & 2025 & Describir y cuantificar los procesos meteorológicos y morfológicos que causan la pérdida de masa en el glaciar Tapado, Chile & Journal of Glaciology \\ \hline

Smart Sensors and Smart Data for Precision Agriculture & Soussi et al. & 2024 & CSintetizar la integración de sensores, Big Data e IA para optimizar la gestión de cultivos y la sostenibilidad & Sensors \\ \hline

Numerical model for the simulation of soil water flow & del Vigo et al. & 2022 & Desarrollar un modelo numérico para simular el flujo de agua en el suelo considerando la absorción de las raíces, aplicado al tomate & Ingeniería del Agua \\ \hline

Glacier and rock glacier changes since the 1950s & Robson et al. & 2022 & Reportar cambios en el área y masa de glaciares en la cuenca de La Laguna, Chile, desde mediados del siglo XX & The Cryosphere \\ \hline  

Explainable Artificial Intelligence in Hydrology & Núñez et al. & 2023 & Interpretar predicciones de caudales por deshielo en una cuenca árida de Chile utilizando técnicas de Inteligencia Artificial Explicable (XAI) & Water \\ \hline

Combining Physical Hydrological Model with Explainable ML & Yang et al. & 2024 & Mejorar la evaluación del balance hídrico en cuencas glaciares mediante un marco híbrido de modelos físicos y aprendizaje automático & Water \\ \hline
\end{longtable}



\subsection{Análisis integrador}

El análisis conjunto de los estudios revisados permite identificar una evolución en los enfoques de modelación aplicados a la gestión hídrica agrícola, los cuales pueden agruparse en tres grandes líneas: modelos biofísicos, modelos de optimización y enfoques basados en comportamiento.

Por una parte, los modelos biofísicos, como AquaCrop, permiten estimar con precisión la relación entre disponibilidad hídrica y rendimiento, proporcionando una base robusta para la evaluación técnica de estrategias de riego. Sin embargo, estos enfoques presentan limitaciones al no incorporar explícitamente la toma de decisiones del agricultor ni las restricciones operativas reales del sistema productivo.

Por otro lado, los enfoques de simulación-optimización integran datos en tiempo real y permiten mejorar la eficiencia del uso del agua, aunque su aplicabilidad en contextos reales se ve restringida por la alta dependencia de datos de alta resolución y monitoreo continuo.

En contraste, los modelos basados en agentes introducen la dimensión del comportamiento humano, evidenciando cómo las decisiones individuales impactan el sistema en su conjunto. No obstante, su implementación requiere una caracterización detallada de los actores, lo que implica una alta complejidad y demanda de información.

Adicionalmente, los estudios sobre sensores IoT evidencian un avance significativo en la disponibilidad de tecnologías para el monitoreo agrícola, aunque su adopción en contextos latinoamericanos sigue siendo limitada debido a barreras económicas, técnicas y de infraestructura.

En este contexto, se observa que ninguno de los enfoques revisados logra, por sí solo, capturar de manera integral la complejidad del sistema agrícola en condiciones de escasez hídrica. En particular, existe una desconexión entre la modelación biofísica del recurso hídrico, la operación productiva y la toma de decisiones de los agricultores.

Por lo tanto, se justifica la necesidad de un enfoque integrador, como la simulación multiparadigma, que permita articular la dinámica del recurso hídrico, la secuencia de procesos productivos y la toma de decisiones en un mismo marco de análisis. Esta integración constituye el principal aporte del presente estudio, diferenciándolo de los enfoques existentes en la literatura.
\chapter{Caracterización y diagnóstico del sistema}

\section{Stakeholders}

Los stakeholders del presente proyecto de investigación corresponden a aquellos actores que podrían estar interesados en sus resultados o verse potencialmente afectados por ellos. En este contexto, se identifican los pequeños y medianos agricultores como principales interesados, dado que los resultados pueden influir en la gestión del riego y la toma de decisiones productivas.

Asimismo, entidades gubernamentales como CEAZA, el Ministerio de Obras Públicas (MOP), la Agencia Nacional de Investigación y Desarrollo (ANID) y la Dirección General de Aguas (DGA) podrían utilizar los resultados para la planificación, regulación y gestión del recurso hídrico. De igual forma, centros de estudio y profesionales como los ingenieros agrónomos podrían verse interesados en la aplicación y desarrollo de herramientas de análisis y simulación en el ámbito agrícola. En conjunto, estos actores representan potenciales usuarios y beneficiarios del conocimiento generado por la investigación.

\section{Oferta Hídrica}

La oferta hídrica constituye uno de los principales factores que condicionan la productividad agrícola en el valle de Pan de Azúcar. A diferencia de otros insumos productivos, su disponibilidad no depende exclusivamente de los procesos hidrológicos naturales, sino también de la interacción entre factores climáticos, infraestructura hidráulica, almacenamiento en embalses, recarga de acuíferos y de las reglas institucionales que regulan la distribución del recurso. \citeref{rivera2025}{Rivera y Donoso (2025)} señalan que la gestión de la sequía en Chile debe entenderse como un problema hidrológico e institucional simultáneamente, donde la disponibilidad efectiva de agua resulta de la interacción entre la oferta natural, la infraestructura de regulación y los mecanismos de administración implementados por las organizaciones de usuarios.

En este contexto, el Código de Aguas establece que el agua constituye un bien nacional de uso público y que el ejercicio de los derechos de aprovechamiento debe compatibilizar la sustentabilidad del recurso, la seguridad hídrica y la satisfacción de las distintas demandas existentes, otorgando prioridad al consumo humano en situaciones de humanas en situaciones de escasez (\citeref{codigo1981}{Código de Aguas}, DFL N° 1.122; Ley N° 21.435, 2022). Como consecuencia, la disponibilidad efectiva de agua para la agricultura no depende únicamente de las condiciones hidrológicas observadas, sino también de las decisiones adoptadas por las organizaciones de usuarios del agua encargadas de su distribución.

En el área de estudio, la principal fuente de abastecimiento superficial corresponde al Canal Bellavista, infraestructura que deriva aguas desde el río Elqui y abastece gran parte de la superficie agrícola del valle de Pan de Azúcar. De acuerdo con \citeref{cnr2015}{la Comisión Nacional de Riego (2015)}, este canal constituye una infraestructura estratégica para el desarrollo agrícola del sector, permitiendo la distribución del recurso mediante un sistema organizado de turnos de riego. Sin embargo, la prolongada sequía que afecta a la Región de Coquimbo ha reducido significativamente la disponibilidad de agua superficial, disminuyendo la superficie efectivamente regada desde aproximadamente 13.000 hectáreas de capacidad potencial a cerca de 4.800 hectáreas actualmente en producción. Esta situación ha obligado a operar con desmarques sucesivos y entregas equivalentes, en numerosos períodos, a entre un 10\% y un 20\% de los derechos nominales de agua, llegando incluso a existir riesgo de suspensión temporal del suministro bajo condiciones hidrológicas extremas.

La disminución sostenida de la oferta superficial ha incrementado la dependencia de fuentes subterráneas, principalmente del acuífero Pan de Azúcar–El Culebrón. No obstante, este acuífero fue declarado área de restricción por la Dirección General de Aguas y permanece cerrado al otorgamiento de nuevos derechos de aprovechamiento desde el año 2000, debido a la creciente presión ejercida sobre el recurso y a la necesidad de resguardar su sustentabilidad. En consecuencia, el agua subterránea constituye una fuente limitada cuya utilización depende principalmente de derechos previamente constituidos y de la capacidad de extracción de los pozos existentes.

\begin{figure}[h]
    \centering
    \includegraphics[width=0.85\textwidth]{fig/DGA2019.png}
    \caption{Condiciones hidrológicas del acuífero El Culebrón 2019}
    \caption*{\footnotesize Fuente: \citeref{dga2019}{Dirección General de Aguas (DGA)}, Boletín Hidrológico N°489.}
    \label{fig:culebron_2019}
\end{figure}

\begin{figure}[h]
    \centering
    \includegraphics[width=0.85\textwidth]{fig/DGA2026.png}
    \caption{Condiciones hidrológicas del acuífero El Culebrón 2026}
    \caption*{\footnotesize Fuente: \citeref{dga2026}{Dirección General de Aguas (DGA)}, Boletín Hidrológico N°573.}
    \label{fig:culebron_2026}
\end{figure}

La comparación entre la Figura~\ref{fig:culebron_2019} y la Figura~\ref{fig:culebron_2026} evidencia la continuidad y profundización de la tendencia descendente observada en los niveles piezométricos del acuífero El Culebrón. Mientras que en el año 2019 el pozo ``El Sauce 2'' ya mostraba una disminución sostenida en sus niveles de agua, para el año 2026 dicha tendencia se intensifica, observándose niveles piezométricos aún menores y una menor capacidad de recuperación entre temporadas.

Estos antecedentes sugieren que la recarga natural del acuífero no logra compensar las extracciones realizadas, reflejando un desequilibrio persistente entre recarga y explotación. \citeref{rivera2025}{Rivera y Donoso (2025)} indican que este comportamiento es consistente con la evolución observada en diversos acuíferos de la zona centro-norte de Chile, donde la combinación entre déficit prolongado de precipitaciones, disminución de la acumulación nival y extracción sostenida de aguas subterráneas ha reducido progresivamente la capacidad de recuperación de los sistemas hidrogeológicos.

En consecuencia, la oferta hídrica disponible para el sistema agrícola del valle de Pan de Azúcar se encuentra condicionada simultáneamente por la disponibilidad superficial proveniente del río Elqui y del Canal Bellavista, por el estado del acuífero Pan de Azúcar–El Culebrón y por las decisiones de gestión adoptadas por las organizaciones de usuarios del agua conforme al marco regulatorio vigente. Esta interacción entre factores hidrológicos, infraestructura hidráulica e institucionalidad constituye el principal origen de la incertidumbre representada posteriormente en el modelo de simulación. 

\section{Demanda hídrica}

Por otra parte, la demanda hídrica en el sector agrícola se encuentra fuertemente influenciada por las condiciones climáticas propias de la región. De acuerdo con \citeref{allen1998}{Allen et al. (1998)}, la evapotranspiración constituye el principal componente de la demanda de agua de los cultivos y depende principalmente de cuatro variables meteorológicas: la radiación solar, la temperatura del aire, la humedad relativa y la velocidad del viento.

En climas semiáridos, como los característicos de la Región de Coquimbo, estas condiciones favorecen tasas elevadas de evapotranspiración. La alta radiación solar y las temperaturas elevadas incrementan la energía disponible para el proceso de evaporación, mientras que la baja humedad relativa y la presencia frecuente de viento aumentan el déficit de presión de vapor de la atmósfera, favoreciendo el transporte de vapor de agua. En particular, la velocidad del viento desempeña un papel fundamental al disminuir la resistencia aerodinámica al intercambio de vapor entre el cultivo y la atmósfera, intensificando las pérdidas de agua desde el sistema suelo--planta--atmósfera (\citeref{allen1998}{Allen et al., 1998}).

A estas condiciones climáticas se suma la variabilidad propia de los cultivos. Según \citeref{allen1998}{Allen et al. (1998)}, la demanda hídrica también depende de características como la especie cultivada, el estado fenológico, la altura del cultivo, el grado de cobertura del suelo y la resistencia aerodinámica del dosel vegetal. Estas características modifican la evapotranspiración del cultivo y generan diferencias significativas en los requerimientos de agua entre especies y entre las distintas etapas de desarrollo.

Desde el punto de vista metodológico, el \citeref{allen1998}{Manual FAO~56} propone estimar la evapotranspiración del cultivo a partir de la evapotranspiración de referencia ($ET_0$) y coeficientes de cultivo ($K_c$). Estos pueden implementarse mediante un coeficiente único o mediante el enfoque de coeficientes duales, el cual separa la evapotranspiración en los procesos de transpiración del cultivo y evaporación del suelo, permitiendo representar con mayor precisión los efectos del desarrollo fenológico, las condiciones de humedad del suelo y las prácticas de manejo del riego. Debido a estas ventajas, el presente trabajo adopta el enfoque de coeficientes duales para la estimación de la demanda hídrica (\citeref{allen1998}{Allen et al., 1998}). 

\section{Irregularidades y condiciones adversas}

Dentro de las principales irregularidades operativas del sistema de riego se identifican los cortes programados del canal destinados a labores de mantención, reparación y ejecución de obras de infraestructura hidráulica. Estas intervenciones son necesarias para garantizar el adecuado funcionamiento del sistema de conducción y distribución del agua, especialmente en canales de riego de larga data como el Canal Bellavista, cuya operación constituye un elemento fundamental para el abastecimiento agrícola del valle de Pan de Azúcar (\citeref{cnr2015}{Comisión Nacional de Riego, 2015}). Si bien estos cortes presentan una duración física aproximada de ocho días, su impacto real sobre los usuarios resulta considerablemente mayor debido al sistema de distribución por turnos implementado en la zona.

En efecto, dado que los agricultores reciben agua en intervalos aproximados de entre siete y nueve días, dependiendo del sector y de la organización de los turnos de riego, la coincidencia de un corte programado con el turno asignado implica la pérdida completa de dicho ciclo de riego. Como consecuencia, los agricultores deben esperar hasta el siguiente turno disponible para acceder nuevamente al recurso, lo que en la práctica puede traducirse en períodos de 15, 18 e incluso hasta 20 días sin disponibilidad de agua para riego.

Los principales afectados corresponden a pequeños y medianos agricultores que dependen exclusivamente del sistema de distribución del canal y que disponen de escasa o nula capacidad de almacenamiento de agua. Esta situación obedece tanto a restricciones económicas como a la limitada incorporación de infraestructura tecnológica en los predios, condición que continúa representando una de las principales brechas identificadas en la agricultura de precisión (\citeref{olivera2025}{Olivera Ruiz \& Pérez Espinoza, 2025}). La ausencia de sistemas de almacenamiento y monitoreo incrementa la vulnerabilidad de estos productores frente a interrupciones del suministro y reduce significativamente su capacidad de adaptación ante escenarios de escasez hídrica.

Asimismo, los desmarques del canal, entendidos como los ajustes periódicos en la asignación efectiva de los derechos de aprovechamiento de agua, constituyen otra fuente relevante de incertidumbre para la planificación agrícola. Estos ajustes son realizados en función de la disponibilidad hídrica proyectada para la temporada y de las condiciones observadas en la cuenca, modificando el porcentaje efectivo de agua que cada agricultor puede utilizar durante el período de riego.

Sin embargo, la periodicidad con que se actualizan estos desmarques no siempre permite responder oportunamente a las variaciones que experimenta el sistema hidrológico durante la temporada agrícola. En consecuencia, los agricultores deben planificar sus actividades productivas bajo condiciones de disponibilidad futura incierta, dificultando decisiones relacionadas con la selección de cultivos, la programación del riego y la asignación de recursos. \citeref{rivera2025}{Rivera y Donoso (2025)} señalan que uno de los principales desafíos de la gestión de la sequía en Chile consiste precisamente en desarrollar herramientas que permitan anticipar escenarios de disponibilidad hídrica y apoyar la toma de decisiones mediante información objetiva, reduciendo la dependencia de estrategias reactivas frente a eventos de escasez.

Bajo estas condiciones, la combinación de cortes operacionales, distribución por turnos, incertidumbre asociada a los desmarques y limitada incorporación de tecnologías de monitoreo configura un entorno altamente complejo para la planificación agrícola. Esta situación refuerza la necesidad de desarrollar herramientas de simulación capaces de integrar la incertidumbre hidrológica, representar distintos escenarios de operación del sistema y apoyar la toma de decisiones de los agricultores antes de la implementación de una estrategia productiva.

\section{Toma de decisiones}

La toma de decisiones en sistemas agrícolas bajo condiciones de escasez hídrica constituye un proceso complejo, debido a que requiere integrar simultáneamente información hidrológica, productiva, económica y climática bajo escenarios caracterizados por una elevada incertidumbre. \citeref{rivera2025}{Rivera y Donoso (2025)} señalan que la gestión de la sequía en Chile ha evolucionado desde un enfoque reactivo hacia la necesidad de desarrollar estrategias preventivas sustentadas en información objetiva y herramientas de apoyo a la decisión, capaces de integrar las múltiples dimensiones que intervienen en la administración del recurso hídrico.

En cuanto a la toma de decisiones a nivel predial, se observa una heterogeneidad significativa entre los distintos tipos de productores. En el caso de los pequeños y medianos agricultores, la planificación agrícola se sustenta principalmente en la experiencia acumulada, el conocimiento empírico y la observación directa de las condiciones del entorno, con una limitada incorporación de herramientas analíticas y tecnologías de monitoreo.

Como consecuencia, las decisiones relacionadas con la selección de cultivos, la programación del riego y la asignación de recursos suelen depender de la percepción individual del agricultor, restringiendo la capacidad de anticipar escenarios futuros y de evaluar sistemáticamente distintas alternativas de manejo.

Por el contrario, los grandes productores disponen de un mayor acceso a tecnologías de monitoreo, sensores, sistemas de información y herramientas de análisis que les permiten integrar continuamente variables hidrológicas, climáticas, productivas y económicas dentro de sus procesos de planificación. Esta disponibilidad de información favorece la adopción de decisiones basadas en evidencia, permitiendo no solo definir qué producir, sino también determinar cuánto producir, cuándo realizar cada labor agrícola y cómo administrar eficientemente los recursos disponibles durante las distintas etapas fenológicas de los cultivos.

Desde el punto de vista metodológico, esta problemática corresponde a un problema de decisión multicriterio, donde la elección de una estrategia productiva depende de la interacción simultánea entre la disponibilidad de agua, los requerimientos hídricos de los cultivos, la estacionalidad de los precios, los costos de producción, la capacidad operativa del predio y la incertidumbre asociada a las condiciones hidrológicas futuras. \citeref{jia2026}{Jia y Zhang (2026)} muestran que este tipo de problemas requiere integrar distintas dimensiones del sistema agrícola para representar adecuadamente las relaciones entre la oferta hídrica, la producción y el desempeño económico.

En este contexto, la simulación constituye una herramienta especialmente adecuada para apoyar la toma de decisiones, ya que permite representar sistemas complejos, incorporar incertidumbre y evaluar múltiples escenarios sin intervenir el sistema real (\citeref{banks2010}{Banks et al., 2010}). De esta forma, es posible analizar anticipadamente distintas estrategias de producción, identificar configuraciones más eficientes y evaluar sus consecuencias técnicas y económicas antes de su implementación, proporcionando información objetiva que complemente la experiencia del agricultor y fortalezca el proceso de planificación agrícola bajo condiciones de escasez hídrica.

\chapter{Opciones de acción}

\section{Evaluación de alternativas de modelación}

En el marco del problema de investigación, se identificaron distintas alternativas de modelación capaces de representar el sistema agrícola bajo condiciones de restricción hídrica. Cada una de estas opciones presenta ventajas específicas, así como limitaciones que deben ser consideradas antes de su implementación.

A continuación, se describen las principales alternativas evaluadas:

\subsection{Dinámica de Sistemas}

La Dinámica de Sistemas es un paradigma de simulación orientado al análisis de sistemas complejos cuyo comportamiento está determinado por relaciones de causalidad, procesos de acumulación y mecanismos de retroalimentación. Su estructura basada en variables de estado, flujos y lazos causales permite representar la evolución temporal de sistemas donde interactúan componentes físicos, ambientales y económicos.

Este enfoque resulta especialmente adecuado para el problema abordado en la presente investigación, ya que la disponibilidad del recurso hídrico depende de procesos continuos como la recarga, el almacenamiento, la extracción y el consumo de agua, cuyos efectos se acumulan y modifican el comportamiento del sistema a lo largo del tiempo. En este contexto, \citeref{jia2026}{Jia y Zhang (2026)} desarrollan un modelo de Dinámica de Sistemas para evaluar la asignación de agua y suelo bajo escenarios de cambio climático, representando explícitamente la interacción entre la disponibilidad hídrica, la humedad del suelo, el rendimiento de los cultivos y los retornos económicos mediante estructuras de retroalimentación. Los autores destacan que este paradigma permite analizar la evolución del sistema y evaluar el efecto de distintas políticas de gestión del agua en horizontes de planificación de largo plazo.

La propuesta de \citeref{jia2026}{Jia y Zhang (2026)} presenta importantes similitudes con la problemática estudiada en esta investigación, ya que ambas buscan representar la interacción entre la disponibilidad de agua, la producción agrícola y las estrategias de gestión frente a escenarios de incertidumbre climática. Asimismo, en ambos casos el interés se centra en comprender cómo las decisiones adoptadas modifican el comportamiento global del sistema a través de relaciones de causa y efecto, más que en representar el comportamiento individual de cada productor.

No obstante, la Dinámica de Sistemas presenta limitaciones cuando es necesario modelar procesos operacionales que ocurren de manera discreta. Si bien permite representar adecuadamente la evolución agregada del sistema y analizar políticas de gestión, no describe con suficiente detalle la secuencia de actividades productivas, la programación del riego, la utilización de recursos o la ejecución de labores agrícolas. De hecho, \citeref{jia2026}{Jia y Zhang (2026)} complementan la Dinámica de Sistemas con técnicas de aprendizaje automático y optimización multiobjetivo para fortalecer el proceso de toma de decisiones, evidenciando que este paradigma suele requerir herramientas adicionales cuando el problema involucra decisiones operacionales específicas. Por esta razón, en la presente investigación se considera pertinente complementar la Dinámica de Sistemas con Simulación de Eventos Discretos, permitiendo representar tanto el comportamiento estratégico del sistema hídrico como la operación detallada del proceso productivo.

\subsection{Modelado Basado en Agentes}

El Modelado Basado en Agentes (Agent-Based Modeling, ABM) constituye una de las principales aproximaciones para representar sistemas socioecológicos complejos, ya que permite modelar explícitamente la heterogeneidad de los actores, sus procesos de toma de decisiones y las interacciones entre ellos y el entorno. En sistemas agrícolas, este paradigma ha sido ampliamente utilizado para analizar cómo las decisiones individuales de los agricultores modifican la disponibilidad de recursos naturales y, simultáneamente, cómo los cambios ambientales retroalimentan el comportamiento de los propios agentes.

Trabajos recientes, como el desarrollado por \citeref{elfartassi2025}{El Fartassi et al. (2025)}, integran un modelo basado en agentes con un modelo biofísico para representar la gestión compartida de recursos hídricos en zonas áridas de Marruecos. En dicho estudio, cada agricultor es caracterizado mediante atributos conductuales asociados a su actitud frente al uso del agua subterránea, capacidad económica, percepción de control sobre la inversión, tipo de tenencia de la tierra e influencia ejercida por otros agricultores, permitiendo simular la adopción de distintas estrategias de riego y el eventual abandono de la actividad agrícola bajo diferentes escenarios climáticos y regulatorios.

Sin embargo, la principal limitación del paradigma radica en la elevada complejidad requerida para caracterizar adecuadamente el comportamiento de los agentes. La definición de reglas de decisión realistas demanda información empírica obtenida mediante encuestas, entrevistas y estudios de comportamiento, además de procesos de calibración y validación que permitan representar de forma confiable la heterogeneidad de la población. Asimismo, la parametrización de variables conductuales y la modelación de las interacciones sociales incrementan considerablemente el esfuerzo de desarrollo y validación del modelo.

Considerando el alcance temporal del presente proyecto, la disponibilidad de información primaria y los recursos destinados a la investigación, la implementación de un modelo basado en agentes no resulta la alternativa más adecuada, debido a la dificultad de caracterizar rigurosamente los procesos de decisión de los agricultores y validar el 
comportamiento emergente del sistema.

\subsection{Simulación de Eventos Discretos}

La Simulación de Eventos Discretos (SED) constituye un paradigma de modelación ampliamente utilizado para representar sistemas cuyo estado evoluciona como consecuencia de la ocurrencia de eventos en instantes específicos del tiempo. De acuerdo con \citeref{banks2010}{Banks et al. (2010)}, un modelo de eventos discretos representa el comportamiento del sistema mediante entidades, recursos, actividades y eventos, donde cada evento modifica el estado del sistema en un instante determinado. Esta característica permite describir con un alto nivel de detalle procesos operacionales, secuencias de actividades, utilización de recursos, tiempos de procesamiento y restricciones propias de sistemas complejos, incorporando además incertidumbre mediante variables estocásticas.

En el ámbito agrícola, \citeref{munyaka2024}{Munyaka et al. (2024)} emplean este paradigma para modelar la cadena de suministro del camote en Zimbabue mediante el software Simul8, representando actividades como la preparación del terreno, cultivo, cosecha, almacenamiento, transporte y comercialización. Además, el modelo incorpora recursos como mano de obra, maquinaria e insumos agrícolas, junto con variables estocásticas relacionadas con el rendimiento de los cultivos, la demanda del mercado y las condiciones climáticas, permitiendo evaluar distintos escenarios, analizar la utilización de recursos e identificar cuellos de botella que afectan el desempeño del sistema.

La propuesta desarrollada por \citeref{munyaka2024}{Munyaka et al. (2024)} presenta importantes similitudes con la problemática abordada en esta investigación, ya que ambas consideran procesos agrícolas compuestos por actividades secuenciales cuya ejecución depende de la disponibilidad de recursos y de decisiones operacionales. En el modelo propuesto, cada cultivo puede representarse como una entidad que transita por distintas etapas del proceso productivo, mientras que recursos como el agua de riego, la mano de obra y la maquinaria condicionan la ejecución de las diferentes labores agrícolas. Esta representación resulta consistente con la estructura conceptual de la Simulación de Eventos Discretos descrita por \citeref{banks2010}{Banks et al. (2010)}, permitiendo modelar explícitamente la interacción entre entidades, recursos y eventos dentro del sistema agrícola.

No obstante, este paradigma presenta limitaciones para representar la dinámica continua del recurso hídrico. Si bien permite modelar adecuadamente la ocurrencia de actividades y eventos operacionales, no describe de forma natural procesos continuos de almacenamiento, disponibilidad, recarga y consumo del agua, ni los mecanismos de retroalimentación que gobiernan la evolución del sistema hidrológico. \citeref{banks2010}{Banks et al. (2010)} señalan que la Simulación de Eventos Discretos resulta especialmente adecuada para sistemas donde predominan eventos discretos y cambios de estado instantáneos, mientras que otros paradigmas ofrecen un mejor desempeño cuando el comportamiento del sistema está dominado por variables continuas. Por esta razón, su utilización como único paradigma resulta insuficiente para responder a los objetivos de la presente investigación, siendo conveniente complementarla con Dinámica de Sistemas para representar el comportamiento continuo de la oferta hídrica y su interacción con el sistema productivo.

\subsection{Simulación Multiparadigma}

La simulación multiparadigma integra distintos enfoques de modelación con el propósito de representar sistemas complejos cuyas componentes presentan comportamientos de naturaleza diferente. En lugar de utilizar un único paradigma para modelar la totalidad del sistema, este enfoque permite seleccionar la técnica de simulación más adecuada para cada subsistema, aprovechando las fortalezas particulares de cada metodología.

La propuesta metodológica desarrollada en esta investigación surge a partir del análisis de trabajos recientes, como el de \citeref{jia2026}{Jia y Zhang (2026)}, quienes abordan el problema de la gestión del agua agrícola descomponiéndolo en diferentes dimensiones, entre ellas la disponibilidad hídrica, la producción agrícola, la evaluación económica y el proceso de toma de decisiones. Cada una de estas dimensiones es modelada mediante herramientas específicas acordes con su naturaleza, evidenciando que la complejidad del sistema puede abordarse de manera más eficiente cuando el problema global se divide en componentes funcionales.

Siguiendo este principio, ampliamente utilizado en ingeniería de sistemas para abordar problemas complejos mediante su descomposición en subproblemas, la presente investigación propone una arquitectura de simulación multiparadigma en la que cada dimensión del sistema agrícola es representada mediante el paradigma más apropiado para su comportamiento. En particular, la dinámica de la oferta hídrica en Chile se modela mediante Dinámica de Sistemas, debido a que corresponde a un fenómeno continuo caracterizado por procesos de acumulación, flujos y mecanismos de retroalimentación asociados a embalses, caudales, derechos de agua, disponibilidad del recurso y restricciones de distribución. Por otra parte, el sistema productivo agrícola se representa mediante Simulación de Eventos Discretos, considerando que las labores agrícolas, la asignación del agua por parte del regante, la utilización de recursos y el desarrollo de los cultivos corresponden a procesos secuenciales gobernados por eventos específicos.

Esta descomposición permite representar de manera independiente la dinámica hidrológica y el comportamiento operacional del sistema agrícola, facilitando posteriormente su integración para generar múltiples escenarios de producción. A partir de estos escenarios es posible identificar patrones de comportamiento y configuraciones productivas que constituyen la base para etapas posteriores de optimización y apoyo a la toma de decisiones. De esta forma, la simulación multiparadigma no solo reduce las limitaciones inherentes al uso de un único paradigma, sino que proporciona una representación más coherente con la naturaleza de cada componente del sistema.

\section{Matriz de selección del paradigma}

Con el objetivo de seleccionar el enfoque metodológico más adecuado, se construyó una matriz de decisión basada en criterios ampliamente utilizados en la evaluación de modelos de simulación, alineados con los requerimientos del problema de investigación.

Los criterios considerados fueron los siguientes:

\begin{itemize}
    \item \textbf{Factibilidad}: Facilidad de implementación técnica del modelo.
    \item \textbf{Viabilidad}: Disponibilidad y acceso a los datos necesarios para su desarrollo.
    \item \textbf{Deseabilidad}: Capacidad del enfoque para representar adecuadamente el sistema real.
    \item \textbf{Complejidad}: Nivel de dificultad conceptual y computacional requerido.
\end{itemize}

A partir de estos criterios, se evaluaron las distintas alternativas metodológicas, obteniéndose la siguiente matriz comparativa:

\begin{table}[H]
\centering
\begin{tabular}{lcccc}
\toprule
\textbf{Paradigma} & \textbf{Factibilidad} & \textbf{Viabilidad} & \textbf{Deseabilidad} & \textbf{Complejidad} \\
\midrule
Dinámica de Sistemas & Alta & Alta & Media & Media \\
Modelado basado en agentes & Baja & Media & Alta & Alta \\
Simulación de eventos discretos & Alta & Alta & Media & Media \\
Simulación multiparadigma & Media & Alta & Alta & Media \\
\bottomrule
\end{tabular}
\caption{Comparación de paradigmas de modelación}
\label{tab:paradigmas}
\end{table}

El análisis de la tabla \ref{tab:paradigmas} evidencia que ningún paradigma individual logra satisfacer simultáneamente los criterios de factibilidad, viabilidad y deseabilidad sin incurrir en altos niveles de complejidad o limitaciones de representación.

En particular, la Dinámica de Sistemas presenta altos niveles de factibilidad y viabilidad, lo que la convierte en una herramienta adecuada para modelar la evolución continua del recurso hídrico. Sin embargo, su deseabilidad es media, dado que no permite representar de manera explícita los procesos operativos ni la ocurrencia de eventos discretos asociados a la gestión agrícola.

Por otro lado, la simulación de eventos discretos destaca por su alta factibilidad y viabilidad, siendo apropiada para modelar procesos productivos y secuencias de eventos. No obstante, al igual que la Dinámica de Sistemas, presenta una deseabilidad media, ya que no logra capturar adecuadamente la dinámica continua del recurso hídrico.

En cuanto al modelado basado en agentes, si bien presenta una alta deseabilidad debido a su capacidad para representar la toma de decisiones y la heterogeneidad de los actores, su baja factibilidad y viabilidad limitan su aplicación en el contexto del estudio, principalmente por la alta demanda de información requerida.

En contraste, la simulación multiparadigma presenta un equilibrio entre los criterios evaluados, destacando por su alta viabilidad y deseabilidad, junto con un nivel de complejidad manejable. Esto sugiere que dicho enfoque permite integrar de manera efectiva las fortalezas de los distintos paradigmas, posicionándose como la alternativa más adecuada según los resultados de la matriz.

\subsection{Justificación}

La selección de la simulación multiparadigma como enfoque metodológico se fundamenta en los resultados obtenidos en el Cuadro \ref{tab:paradigmas}, donde se observa que este paradigma es el único que logra un desempeño equilibrado en todos los criterios de evaluación, sin presentar debilidades críticas en ninguno de ellos.

A diferencia de los enfoques individuales, que presentan limitaciones estructurales en su capacidad de representación, la simulación multiparadigma permite abordar el problema de investigación desde una perspectiva integradora. En particular, combina la capacidad de la Dinámica de Sistemas para modelar la evolución continua del recurso hídrico con la fortaleza de la simulación de eventos discretos para representar procesos operativos y decisiones agrícolas.

Esta integración resulta especialmente relevante en el contexto del estudio, donde el sistema presenta una naturaleza híbrida, caracterizada por la interacción entre procesos físicos continuos y decisiones discretas de gestión. En este sentido, el enfoque multiparadigma no solo mejora la capacidad descriptiva del modelo, sino que también permite capturar las interdependencias entre la disponibilidad de agua y las acciones productivas.

En consecuencia, la adopción de este enfoque metodológico permite desarrollar un modelo más coherente con la realidad del sistema, fortaleciendo su validez analítica y su utilidad como herramienta para la toma de decisiones en escenarios de escasez hídrica. 

\subsection{Supuestos}

Con el fin de hacer operativa la modelación del sistema, se establecen los siguientes supuestos, los cuales permiten simplificar ciertos procesos sin comprometer la capacidad del modelo para representar la dinámica general del sistema:

\begin{enumerate}

    \item \textbf{Se consideran sistemas de riego sin fallas.}  
    Este supuesto permite aislar el efecto de la disponibilidad hídrica y las estrategias de riego sobre el sistema productivo. La inclusión de fallas operativas introduciría una fuente adicional de incertidumbre que escapa al foco del estudio. No obstante, se reconoce que estas pueden existir en la práctica y podrían ser incorporadas en futuras extensiones del modelo.

    \item \textbf{El crecimiento del coeficiente de cultivo entre fases es lineal.}  
    Se asume una evolución lineal del coeficiente de cultivo (Kc) entre etapas fenológicas como una aproximación simplificada de su comportamiento real. Esta decisión se fundamenta en la necesidad de interpolar valores entre etapas discretas definidas en la literatura (FAO 56), permitiendo una implementación continua en el modelo.

    \item \textbf{El cultivo se encuentra al aire libre.}  
    Este supuesto permite considerar directamente la influencia de variables climáticas sobre la evapotranspiración. Se excluyen sistemas protegidos, como invernaderos, debido a que requieren una modelación específica. 
    
    \item \textbf{Toda la producción obtenida es comercializada.}  
    Se asume que la totalidad de la producción cosechada encuentra un mercado para su venta, por lo que no se consideran pérdidas asociadas a problemas de comercialización, almacenamiento o demanda. En consecuencia, las diferencias en los ingresos entre alternativas se deben exclusivamente a la cantidad producida, la calidad alcanzada y el precio de venta correspondiente al mes de cosecha. Este supuesto permite centrar el análisis en el efecto de la disponibilidad hídrica y las decisiones productivas sobre el desempeño económico del sistema.

\end{enumerate}

\section{Evaluación de alternativas de implementación}

Una vez seleccionado el paradigma de \textbf{simulación multiparadigma}, corresponde definir la \textbf{plataforma tecnológica} sobre la cual materializarlo. De forma análoga a la selección del paradigma, se evaluaron distintas alternativas de implementación considerando su capacidad para integrar los tres modelos requeridos (predictivo, dinámica de sistemas y eventos discretos), su flexibilidad, su costo y su viabilidad para desplegarse como un Sistema de Apoyo a la Decisión (DSS) accesible para el agricultor. Se contrastaron tres alternativas: \textbf{AnyLogic}, \textbf{AquaCrop} y \textbf{Python}.

\subsection{AnyLogic}

AnyLogic es un software comercial especializado en simulación que soporta de forma nativa la dinámica de sistemas, la simulación de eventos discretos y el modelado basado en agentes dentro de un mismo entorno visual. Su principal fortaleza es la integración multiparadigma y su interfaz gráfica.

No obstante, presenta limitaciones relevantes para este proyecto: requiere una licencia comercial de alto costo, ofrece menor flexibilidad para implementar lógica personalizada \textemdash como el algoritmo de selección combinatoria de cultivos o el modelo predictivo del desmarque\textemdash\ y dificulta el despliegue del resultado como una aplicación web de libre acceso para los usuarios finales.

\subsection{AquaCrop}

AquaCrop es un modelo biofísico desarrollado por la FAO para simular la respuesta del rendimiento de los cultivos a la disponibilidad de agua. Se encuentra ampliamente validado y es robusto en la relación agua--rendimiento.

Sin embargo, es un modelo cerrado y de propósito específico: no permite representar la oferta hídrica del canal (turnos, desmarque, pérdidas y almacenamiento), no incorpora el paradigma de eventos discretos ni la lógica de decisión económica, y no contempla la construcción de un DSS. Por ello, resulta insuficiente para abordar de manera integral el problema planteado.

\subsection{Python}

Python es un lenguaje de programación de propósito general y de código abierto que dispone de un ecosistema de librerías capaz de cubrir los tres paradigmas requeridos dentro de un mismo entorno. Permite implementar la dinámica de sistemas, la simulación de eventos discretos, el modelo predictivo, el balance hídrico FAO-56, la optimización combinatoria y el sistema de apoyo a la decisión, sin costos de licencia y con total flexibilidad. Además, favorece la trazabilidad, la reproducibilidad y el despliegue del modelo como una herramienta web accesible.

\subsection{Matriz de selección de la plataforma}

Con los mismos criterios utilizados para el paradigma, se comparan las tres alternativas de implementación:

\begin{table}[H]
\centering
\begin{tabular}{lcccc}
\toprule
\textbf{Plataforma} & \textbf{Integración multiparadigma} & \textbf{Flexibilidad} & \textbf{Costo} & \textbf{Despliegue como DSS} \\
\midrule
AnyLogic & Alta & Media & Alto & Baja \\
AquaCrop & Baja & Baja & Bajo & Baja \\
Python   & Alta & Alta & Bajo & Alta \\
\bottomrule
\end{tabular}
\caption{Comparación de plataformas de implementación}
\label{tab:plataformas}
\end{table}

\subsection{Justificación de la plataforma seleccionada}

La comparación de la Tabla \ref{tab:plataformas} evidencia que Python es la única alternativa que satisface simultáneamente la integración de los tres paradigmas, la flexibilidad para implementar la lógica específica del modelo y la posibilidad de desplegar un DSS accesible, todo ello sin costos de licencia. AnyLogic, si bien potente en la integración multiparadigma, resulta costoso y poco flexible para el despliegue requerido; AquaCrop, aunque robusto, es demasiado específico y no cubre la oferta hídrica ni la toma de decisiones. En consecuencia, se seleccionó \textbf{Python} como plataforma de implementación.

\subsection{Librerías utilizadas por paradigma}

Dentro del ecosistema de Python, cada componente del modelo se implementó con librerías especializadas, seleccionadas por su madurez y adecuación a cada paradigma. La Tabla \ref{tab:librerias} resume las alternativas disponibles y la finalmente adoptada en cada caso.

\begin{table}[H]
\centering
\small
\begin{tabular}{p{4.5cm} p{4.5cm} p{3.5cm}}
\toprule
\textbf{Componente / Paradigma} & \textbf{Librerías disponibles} & \textbf{Seleccionada} \\
\midrule
Modelo predictivo del desmarque (regresión) & statsmodels, scikit-learn, scipy & statsmodels y scikit-learn \\
Dinámica de sistemas (oferta hídrica) & PySD, BPTK-Py & PySD \\
Eventos discretos (proceso productivo) & SimPy, Salabim & SimPy \\
Manejo e integración de datos & pandas, numpy, openpyxl, requests & Todas \\
Visualización y DSS & Dash + Plotly, Streamlit, matplotlib & Dash + Plotly \\
\bottomrule
\end{tabular}
\caption{Librerías de Python evaluadas y seleccionadas por paradigma}
\label{tab:librerias}
\end{table}

En síntesis, se emplea \textbf{statsmodels} (junto con \textbf{scikit-learn} y \textbf{scipy}) para el modelo predictivo del desmarque, \textbf{PySD} para la dinámica de sistemas de la oferta hídrica, \textbf{SimPy} para la simulación de eventos discretos del proceso productivo, la pila \textbf{pandas}/\textbf{numpy} para el manejo de datos y \textbf{Dash} con \textbf{Plotly} para la construcción del sistema de apoyo a la decisión.

\section{Alcances y limitaciones del método}

El principal aporte metodológico de esta investigación consiste en la integración interdisciplinaria de Sistemas de Información Geográfica (SIG), la geometría y planimetría oficial de la red de canales, parámetros hidráulicos empíricos de pérdidas por conducción y un enfoque de simulación multiparadigma. Esta integración permite estructurar un Sistema de Soporte a las Decisiones (DSS) agrícolas que evalúa de forma conjunta la variabilidad hidrológica exógena, la restricción física espacial de la infraestructura hidráulica y la programación operacional y económica predial (Brailsford et al., 2019; Law, 2015).

Esta integración permite evaluar de forma conjunta la disponibilidad de agua, el desarrollo fenológico de los cultivos, la calidad de la producción y el desempeño económico de distintas alternativas agrícolas bajo escenarios de incertidumbre hídrica, proporcionando una herramienta de apoyo a la toma de decisiones para pequeños y medianos agricultores. Asimismo, la interacción entre los distintos subsistemas facilita la representación de relaciones de dependencia entre la disponibilidad del recurso hídrico y las decisiones productivas, permitiendo analizar de manera integrada los efectos de distintos escenarios de gestión del agua (\citeref{jia2026}{Jia \& Zhang, 2026}).

Entre los principales alcances del método destaca la incorporación explícita de la incertidumbre asociada al porcentaje de desmarque mediante un modelo predictivo, cuya variabilidad es propagada hacia la simulación de la oferta hídrica. Este enfoque permite representar de mejor manera la incertidumbre inherente a la disponibilidad de agua y evaluar su impacto sobre el desempeño técnico y económico de las distintas alternativas agrícolas.

Asimismo, la generación exhaustiva de combinaciones de cultivo, junto con la representación de la estacionalidad de los precios, la calidad comercial de la producción y las restricciones presupuestarias, permiten evaluar integralmente distintas estrategias de producción dentro de un único marco de simulación. Esta capacidad de explorar un amplio conjunto de escenarios constituye una de las principales fortalezas de la simulación como herramienta de apoyo a la toma de decisiones, al permitir identificar patrones de comportamiento y alternativas de manejo difíciles de analizar mediante procedimientos analíticos convencionales (Law, 2015; Brailsford et al., 2019).

No obstante, el modelo presenta algunas limitaciones derivadas de la disponibilidad de información y de las decisiones metodológicas adoptadas durante su desarrollo. La estimación de la evapotranspiración se realiza utilizando el procedimiento descrito en el \citeref{allen1998}{Manual FAO~56}, el cual constituye una metodología ampliamente validada y utilizada en estudios de gestión del riego (\citeref{allen1998}{Allen et al., 1998}).

Sin embargo, durante las reuniones de levantamiento de requerimientos realizadas con la empresa PTI Hortícola se identificó el potencial de metodologías basadas en sensores remotos para complementar la estimación de las necesidades hídricas de los cultivos. En particular, índices de vegetación como el \textit{Normalized Difference Vegetation Index} (NDVI) han demostrado ser herramientas eficaces para el monitoreo del estado fisiológico de los cultivos, la caracterización del vigor vegetal y el apoyo a la estimación de parámetros agronómicos relacionados con los requerimientos hídricos, permitiendo además ampliar la cobertura espacial y la frecuencia de actualización de la información disponible (\citeref{dingre2020}{Dingre \& Gorantiwar, 2020}; ...). De manera complementaria, \citeref{munyaka2024}{Munyaka et al. (2024)} integran información proveniente de imágenes Landsat, MODIS e indicadores climáticos derivados de sensores remotos dentro de un modelo de simulación agrícola para representar el efecto de las condiciones ambientales sobre la productividad y la resiliencia de distintas variedades de cultivo, evidenciando el creciente uso de la teledetección como fuente de información para fortalecer modelos de simulación y herramientas de apoyo a la toma de decisiones. Asimismo, ...

En este contexto, la incorporación de un módulo de estimación de la evapotranspiración apoyado en imágenes satelitales constituye una línea de investigación futura e incluso una oportunidad para el desarrollo de un \textit{spin-off} orientado al monitoreo agrícola y a la actualización automática de parámetros del modelo.

Otra limitación corresponde a la representación del comportamiento comercial de los cultivos. En la presente investigación se asume que toda la producción obtenida logra ser comercializada, diferenciando únicamente el valor económico capturado mediante la calidad alcanzada y la estacionalidad de los precios históricos.

Si bien este supuesto permite centrar el análisis en la planificación agrícola bajo condiciones de incertidumbre hídrica, una futura extensión del modelo podría incorporar un indicador de comercialización construido a partir de información histórica de mercados, demanda estacional, frecuencia de transacciones y facilidad de venta de cada hortaliza. La incorporación de este criterio permitiría complementar la evaluación económica con una estimación explícita del potencial comercial de cada alternativa, fortaleciendo aún más el proceso de apoyo a la toma de decisiones.

Asimismo, el modelo podría ampliarse mediante la incorporación de un subsistema destinado a representar la capacidad operativa del predio agrícola, considerando variables como disponibilidad de mano de obra, utilización de maquinaria, tiempos de preparación del terreno, trasplante, labores culturales y cosecha. Esto permitiría representar de manera más realista las restricciones operacionales asociadas a la gestión simultánea de múltiples cultivos y mejorar la programación de las actividades productivas, aspecto que ha demostrado ser adecuadamente representado mediante Simulación de Eventos Discretos en sistemas agrícolas (\citeref{munyaka2024}{Munyaka et al., 2024}).

Una limitación inherente a la componente geoespacial del modelo de oferta hidráulica corresponde a que la red y geometría de conducción implementadas representan exclusivamente el levantamiento catastral oficial disponible al momento de este estudio. Por consiguiente, los coeficientes de pérdida por conducción adoptados corresponden a los aforos experimentales documentados en el informe técnico de \citeref{arcadis2015}{ARCADIS (2015)} para el Canal Bellavista. Cualquier modificación física futura en la sección transversal del canal, labores de revestimiento de tramos de tierra a hormigón, o instalación de nuevas compuertas automatizadas de derivación requerirá una recalibración directa de los parámetros de eficiencia hidráulica ($\eta(L)$) y la actualización de las capas del archivo KML para mantener la representatividad física de la simulación.

Finalmente, la precisión del modelo puede incrementarse mediante campañas experimentales desarrolladas en la Región de Coquimbo, orientadas a recopilar información local sobre humedad del suelo, calendarios de riego, desarrollo fenológico, rendimiento y calidad de los cultivos. Asimismo, la incorporación de sensores de terreno e información proveniente de plataformas de observación satelital permitiría recalibrar periódicamente los parámetros agronómicos del modelo, disminuyendo la dependencia de antecedentes provenientes de fichas técnicas y fortaleciendo progresivamente su capacidad predictiva para condiciones reales de producción (...).

En consecuencia, los resultados obtenidos deben interpretarse como una herramienta de apoyo a la toma de decisiones y no como una representación exacta del sistema agrícola. Sin embargo, la arquitectura multiparadigma propuesta permite integrar procesos hidrológicos, operacionales y económicos dentro de un único marco metodológico, posibilitando el análisis de múltiples escenarios de disponibilidad hídrica, la exploración sistemática de distintas combinaciones de cultivo y la generación de información objetiva que complemente la experiencia del agricultor. Estas características constituyen un aporte significativo para la planificación agrícola bajo condiciones de incertidumbre y proporcionan una base metodológica sólida para futuras extensiones del modelo, incorporando nuevas fuentes de información, técnicas de optimización y herramientas de monitoreo agrícola (...).

\chapter{Materiales y metodos}

\section{Área de estudio o contexto}

El área de estudio se localiza en una zona representativa de clima semiárido, caracterizada por una marcada variabilidad interanual en las precipitaciones y una alta dependencia de los recursos hídricos superficiales y subterráneos. 

La altitud del sitio se sitúa aproximadamente 500 metros sobre el nivel del mar, lo que influye directamente en las condiciones térmicas y en los procesos de evapotranspiración. Las temperaturas presentan una amplitud moderada, con veranos cálidos e inviernos suaves, mientras que las precipitaciones se concentran principalmente en los meses invernales.

La zona se encuentra bajo la influencia de factores climáticos regionales que determinan la disponibilidad hídrica y la estacionalidad de los cultivos. 

\section{Nomenclatura}

A continuación se define la notación empleada a lo largo del modelo, agrupada según su rol dentro del proceso. Estos símbolos se utilizan de manera consistente en las secciones posteriores, tanto en el \textbf{balance hídrico} como en la \textbf{estructura general del proceso} y en la formulación de los \textbf{indicadores} de desempeño.

\vspace{0.2cm}

\subsection*{Índices y dimensión temporal}

\begin{itemize}
    \item $t$: día de simulación (paso temporal diario).
    \item $i$: índice de escenario de desmarque, con $i \in \{-2,-1,0,1,2\}$.
    \item $n$: número de cultivos factibles para la fecha de siembra.
    \item $k$: número de particiones en que se divide el terreno.
\end{itemize}

\subsection*{Oferta hídrica y desmarque}

\begin{itemize}
    \item $D$: porcentaje de desmarque (fracción de los derechos de agua utilizable).
    \item $\hat{D}$: desmarque estimado por el modelo de regresión.
    \item $d_0$: valor base de desmarque definido por el regante.
    \item $\Delta d$: salto entre escenarios de desmarque.
    \item $I$: indicador sintético de una variable hidrológica por temporada.
    \item $\bar{x},\ x_{\min},\ x_{\max}$: valor promedio, mínimo y máximo de dicha variable en la temporada.
\end{itemize}

\subsection*{Demanda hídrica (FAO-56)}

\begin{itemize}
    \item $ET_0$: evapotranspiración de referencia [mm/día].
    \item $ET_c$: evapotranspiración del cultivo [mm/día].
    \item $K_{cb}$: coeficiente basal del cultivo.
    \item $K_e$: coeficiente de evaporación del suelo.
    \item $K_c$: coeficiente de cultivo, con $K_c = K_{cb} + K_e$.
\end{itemize}

\subsection*{Balance hídrico y suelo}

\begin{itemize}
    \item $H_t$: humedad del suelo en el día $t$.
    \item $P_t$: precipitación efectiva.
    \item $R_t$: lámina de riego aplicada.
    \item $D_t$: déficit hídrico respecto de la capacidad de campo.
    \item $CC$: capacidad de campo (límite superior $H_{\max}$).
    \item $PMP$: punto de marchitez permanente (límite inferior $H_{\min}$).
    \item $AFA$: agua fácilmente aprovechable.
    \item $ATA$: agua total aprovechable.
\end{itemize}

\subsection*{Indicadores económicos y de calidad}

\begin{itemize}
    \item $P_1,\ P_2$: proporción de producción de primera y segunda calidad, respectivamente.
    \item $FCV$: factor de captación del valor, $FCV = P_1 + 0.60\,P_2$.
    \item $IC$: indicador de calidad compuesta, $IC = 0.60\,P_1 + 0.40\,P_2$.
    \item $Ingreso_{neto}$: ingreso neto, $Ingreso_{neto} = Ingreso - Costo$.
    \item $Score$: indicador de desempeño global, $Score = Ingreso_{neto}\cdot IC/100$.
    \item $\eta_{tec}$: factor de eficiencia por tecnificación del riego (goteo = 0{,}95; aspersión = 0{,}80; tradicional = 0{,}60).
\end{itemize}
\section{Fuentes de datos}

La investigación utiliza múltiples fuentes de información clave:

\begin{itemize}
    \item \textbf{Estación meteorológica CEAZA}: Proporciona información en tiempo real e histórica de múltiples variables de interés para la investigación. 
    
    \item \textbf{Boletines Hidrométricos DGA}: Permiten caracterizar la disponibilidad y variabilidad del recurso hédrico en la zona.
    
    \item \textbf{Entrevistas a agricultores}: Recopilan información sobre prácticas de manejo, toma de decisiones y percepciones frente a escenarios de escasez hédrica. 
    
    \item \textbf{Web Oficina de Estudios y Políticas Agrarias}: Proporciona información historica de los precios para la selección de cultivos.
    
\end{itemize}



\section{Herramientas utilizadas}

El desarrollo del proyecto se sustenta principalmente en el lenguaje \textbf{Python}, empleado como base para la construcción del modelo multiparadigma, el modelo predictivo y los distintos productos de visualización. A continuación se detallan las herramientas y librerías utilizadas, agrupadas según su función dentro del proyecto.

\vspace{0.2cm}

\subsection*{Modelo predictivo del desmarque}

Para la estimación del porcentaje de desmarque mediante regresión lineal múltiple se emplearon las siguientes librerías:

\begin{itemize}
    \item \textbf{pandas} y \textbf{numpy}: manipulación de las series históricas hidrológicas y meteorológicas, construcción de la matriz de predictores y cálculo numérico.
    \item \textbf{statsmodels}: ajuste del modelo por mínimos cuadrados ordinarios (OLS) y verificación de los supuestos estadísticos mediante las pruebas de Jarque--Bera, Breusch--Pagan y Durbin--Watson.
    \item \textbf{scikit-learn}: apoyo en los procedimientos de selección de variables (\textit{forward}, \textit{backward} y \textit{stepwise}) y en el cálculo de métricas de error ($RMSE$ y $R^2$).
    \item \textbf{scipy}: funciones estadísticas complementarias para el análisis de los residuos.
\end{itemize}

\subsection*{Simulación multiparadigma}

\begin{itemize}
    \item \textbf{PySD}: motor de dinámica de sistemas empleado para modelar la evolución temporal de la oferta hídrica (stocks y flujos del recurso superficial y subterráneo).
    \item \textbf{simpy}: motor de simulación de eventos discretos que representa el ciclo productivo de los cultivos y la operación del sistema de riego.
    \item \textbf{pandas} y \textbf{numpy}: gestión del balance hídrico diario, manejo de los datos de entrada y consolidación de resultados.
    \item \textbf{openpyxl}: lectura y escritura de archivos Excel con parámetros y resultados del modelo.
    \item \textbf{requests}: extracción automatizada de datos climáticos e hidrológicos desde las plataformas de CEAZA y la DGA.
\end{itemize}

\subsection*{Procesamiento y preparación de datos geoespaciales}

Para la preparación, georreferenciación y procesamiento topológico de la red hidráulica se emplearon las siguientes herramientas y librerías:

\begin{itemize}
    \item \textbf{AutoCAD (Autodesk)}: Edición y depuración de la planimetría catastral oficial del Canal Bellavista suministrada en formato de dibujo vectorial (DWG).
    \item \textbf{Global Mapper}: Conversión de archivos CAD vectoriales al formato geoespacial Keyhole Markup Language (KML), proyectando y georreferenciando las coordenadas del trazado del canal principal y sus bifurcaciones secundarias en el valle de Pan de Azúcar.
    \item \textbf{GeoPandas}, \textbf{Shapely} y \textbf{PyOGRIO} (en Python): Reconstrucción topológica de la red de canales en un grafo de flujo, proyección de coordenadas espaciales a la proyección local UTM (EPSG:32719) y cálculo de distancias hidráulicas lineales mínimas sobre la red desde la bocatoma hasta la entrega predial de cada regante.
\end{itemize}

\subsection*{Visualización y sistema de apoyo a la decisión (DSS)}

\begin{itemize}
    \item \textbf{Dash} y \textbf{dash\_cytoscape}: construcción de la interfaz web interactiva del DSS y de los diagramas del modelo.
    \item \textbf{Plotly}: gráficos interactivos de los reportes (escenarios, calidad y presupuesto).
    \item \textbf{matplotlib}: generación de gráficos estáticos para el análisis y los reportes.
    \item \textbf{Flask}, \textbf{flask-cors} y \textbf{Werkzeug}: backend del servidor que expone las animaciones y gestiona la comunicación con el \textit{frontend}.
\end{itemize}

\subsection*{Creación de material divulgativo}

\begin{itemize}
    \item \textbf{Manim}: generación programática de animaciones explicativas del modelo (desmarque, dinámica de sistemas y flujo general del proceso).
    \item \textbf{Google Earth Studio}: Creación de escenas de recorrido satelital.  
    \item \textbf{Sony Vegas}: edición y post-producción del material audiovisual de divulgación. 
    
\end{itemize}

\subsection*{Herramientas organizacionales y de apoyo}

\begin{itemize}
    \item \textbf{Trello}: organización del trabajo, asignación de tareas y seguimiento del progreso bajo el enfoque Scrum.
    \item \textbf{Google Earth}: caracterización de predios y mapeo de la ruta del agua mediante imágenes satelitales.
\end{itemize}


\section{Metodología de desarrollo de la solución}

Antes de detallar la formulación del modelo, esta sección presenta el \textbf{marco de trabajo} bajo el cual fue construido. El desarrollo se abordó de manera iterativa, articulando la gestión del proyecto "\textit{qué} se hizo y en qué orden" con el proceso técnico de construcción de la simulación "\textit{cómo} se formalizó". A partir de este marco general, las secciones siguientes avanzan progresivamente hacia la estructura del proceso, su formulación matemática, su implementación y su validación.

\vspace{0.2cm}

\subsection{Enfoque de trabajo: Scrum Thinking}

Se definió para el desarrollo del proyecto el uso de Scrum Thinking, metodología iterativa con fuerte enfoque en el usuario, con el fin de alcanzar la innovación desde el diseño, integrando las necesidades de las personas, las posibilidades de la tecnología y los requerimientos del negocio para alcanzar el éxito.

Se basa en la integración de las metodologías Scrum y Design Thinking, de tal modo que el inicio del proceso se fundamenta en esta última y su enfoque en la empatización con el usuario, con el fin de recolectar todos los puntos importantes para definir la problemática e idear una solución acorde a la misma. Para el desarrollo de la solución, se toma desde Scrum el uso de sprints, los cuales constan de etapas no mayores a 2 a 4 semanas para construir funcionalidades. De esta forma, se realizan pruebas y validaciones en cortos períodos de tiempo, fomentando la adaptabilidad y la corrección temprana, finalizando nuevamente en la empatización para reiniciar el proceso e identificar nuevos aspectos de mejora.

\begin{figure}[h]
    \centering
    \includegraphics[width=0.85\textwidth]{fig/Scrum thinking.png}
    \caption{Integración de Design Thinking y Scrum en el desarrollo iterativo}
    \caption*{\footnotesize Fuente: Adaptado de Coach Takeshi, \textit{Design Thinking + Scrum}.}
    \label{fig:scrum_thinking}
\end{figure}

Su aplicación dentro del proyecto se estructura en base a las fases de Scrum Thinking. Para la etapa inicial, se considera la realización de entrevistas a agricultores, con el fin de comprender la problemática de la escasez hídrica y las prácticas de riego en el contexto real del sector, identificando los puntos más relevantes a abordar junto con las variables del modelo a considerar.
Posteriormente, para el desarrollo y construcción de los modelos de simulación, se organiza el trabajo en sprints iterativos, donde, en pos de la mejora continua, se busca acercarse al final de cada ciclo al resultado óptimo a validar en la etapa final. Esto se realiza mediante la comparación de los resultados con datos históricos, junto con la experiencia de los agricultores, asegurando su coherencia y utilidad para la toma de decisiones. En caso de ser necesario, se reevaluarán funcionalidades en base a la retroalimentación y pruebas de funcionamiento.






\section{Arquitectura general del modelo de simulación}

El desarrollo de la solución sigue un enfoque metodológico sistemático de simulación estructurado en etapas iterativas (Conceptualización, Recolección de datos, Construcción, Verificación/Validación y Diseño experimental), como se ilustra en el flujo conceptual de la Figura~\ref{fig:proceso_simulacion}.

\begin{figure}[h]
    \centering
    \includegraphics[width=0.7\textwidth]{fig/Metodologia simulación.png}
    \caption{Etapas metodológicas del desarrollo de la simulación}
    \caption*{\footnotesize Fuente: Adaptado de Banks et al., \textit{Discrete-Event System Simulation}.}
    \label{fig:proceso_simulacion}
\end{figure}

El modelo desarrollado se concibe como un \textbf{proceso secuencial} que, bajo un enfoque de \textbf{simulación multiparadigma}, articula tres grandes bloques: la representación de la \textbf{oferta hídrica}, la \textbf{simulación del proceso productivo} y la \textbf{selección de la estrategia de cultivo} más conveniente. En la base del proceso, la \textbf{dinámica de sistemas} representa la evolución temporal de la oferta hídrica disponible para el riego, incorporando los principales procesos asociados a la distribución del recurso superficial y subterráneo. No obstante, uno de los parámetros que condiciona este comportamiento es el \textbf{porcentaje de desmarque}, cuya incertidumbre depende de las condiciones hidrológicas de la cuenca. Para representar dicha incertidumbre de forma objetiva, se desarrolló un \textbf{modelo predictivo} que estima el porcentaje de desmarque a partir de información hidrológica y climática histórica, empleando dicha estimación y su error de predicción como referencia para construir distintos \textbf{escenarios de disponibilidad hídrica}. A partir de estos escenarios, el modelo simula el \textbf{proceso productivo agrícola} y evalúa diferentes alternativas de cultivo bajo restricciones hídricas, técnicas y económicas, finalizando con la selección de la estrategia productiva de mejor desempeño.

\section{Construcción del modelo de oferta hídrica}

El modelo de oferta de agua estima la disponibilidad diaria del recurso a nivel de predio. Su construcción requiere el acoplamiento de la dinámica hidrológica de la cuenca, la predicción de los desmarques institucionales y la restricción física espacial de la red de canales.

\subsection{Modelación hidrológica de la oferta}

La oferta hídrica es representada mediante un modelo de \textbf{dinámica de sistemas} que simula la disponibilidad diaria de agua para riego durante toda la temporada agrícola. El modelo incorpora las principales reglas de operación del sistema de distribución, permitiendo representar la evolución temporal de la oferta superficial y subterránea disponible para el agricultor.

\vspace{0.2cm}

Dentro del modelo se consideran, entre otros, los siguientes \textbf{parámetros operativos}:

\begin{itemize}
    \item \textbf{Número de acciones de agua} y su volumen asociado.
    \item \textbf{Porcentaje de desmarque} y \textbf{fecha de cambio} del mismo.
    \item \textbf{Ubicación del predio} (coordenadas geográficas del regante) para la determinación de pérdidas físicas en la red.
    \item \textbf{Ciclo de turnos} de distribución del canal.
    \item \textbf{Pérdidas} por conducción asociadas a la longitud hidráulica recorrida.
    \item \textbf{Eventos de recarga} de aguas subterráneas.
\end{itemize}

A partir de estos elementos y de la cota máxima hidráulica ($Q_{max}$) determinada por la localización predial en la red de canales, el modelo calcula diariamente la \textbf{oferta efectiva} de agua disponible para riego, generando como salida un calendario de oferta el cual será el dato de entrada para el modelo de simulación agrícola.

Cabe destacar que el impacto de este calendario de oferta sobre la planificación predial varía según la ventana temporal del ciclo productivo. Para aquellos cultivos cuyo desarrollo ocurre íntegramente dentro de un intervalo de desmarque ya decretado (por ejemplo, cosechas de verano desarrolladas entre enero y marzo), la dotación asignada es conocida y no presenta incertidumbre operativa. Sin embargo, para aquellos cultivos cuyos ciclos se extienden o inician en fechas cercanas al mes de septiembre —período en que la Junta de Vigilancia del Río Elqui actualiza anualmente la tasa de desmarque— existe un nivel significativo de incertidumbre respecto del caudal que efectivamente recibirán durante el resto de la temporada. 

La Figura~\ref{fig:evolucion_desmarque_junta} ilustra la evolución histórica de las tasas de desmarque decretadas por la Junta de Vigilancia del Río Elqui (\citeref{jvre2025}{Junta de Vigilancia del Río Elqui, 2025}), evidenciando fluctuaciones significativas e imprevistas entre distintas temporadas hidrológicas. Esta oscilación histórica fundamenta la necesidad metodológica de contar con un modelo de predicción capaz de anticipar el porcentaje de desmarque y alimentar los escenarios de simulación para la toma de decisiones agrícolas.

\begin{figure}[H]
    \centering
    \includegraphics[width=0.85\textwidth]{fig/desmarque.png}
    \caption{Evolución histórica del porcentaje de desmarque en la cuenca del Río Elqui.}
    \caption*{\footnotesize Fuente: \citeref{jvre2025}{Junta de Vigilancia del Río Elqui (2025)}.}
    \label{fig:evolucion_desmarque_junta}
\end{figure}

\vspace{0.2cm}

\subsection{Modelo predictivo del desmarque}

Uno de los parámetros de \textbf{mayor influencia} en la estimación de la oferta hídrica corresponde al \textbf{porcentaje de desmarque} ($D$), ya que determina la fracción efectiva de los derechos de aprovechamiento de aguas superficiales que puede utilizarse durante la temporada agrícola. Como parte del levantamiento de antecedentes de la investigación, se realizaron reuniones técnicas con profesionales de la Junta de Vigilancia del Río Elqui, entre ellos Alex Cortés, en las cuales se detalló el procedimiento mediante el cual la organización determina anualmente esta tasa. Estas sesiones permitieron comprender que la fijación del desmarque no responde a una regla determinística ni a un umbral fijo para una única variable, sino a la evaluación conjunta de diversos indicadores hidrológicos y meteorológicos de la cuenca, constituyendo esta comprensión la base fundamental para la identificación y selección inicial de las variables explicativas consideradas en el modelo predictivo.

Específicamente, el porcentaje de desmarque es actualizado anualmente durante el mes de septiembre por la Junta de Vigilancia a partir de las condiciones hidrológicas observadas en la temporada procedente y de las proyecciones de disponibilidad de agua para el período siguiente. A partir del levantamiento realizado con la organización, se identificó que las principales variables que intervienen en esta evaluación multicriterio corresponden al volumen almacenado en el embalse Puclaro, la acumulación nival en la cordillera, los caudales observados en el río Elqui y las condiciones meteorológicas locales. La interacción entre estos factores determina finalmente la asignación hídrica aprobada, incorporando un nivel importante de incertidumbre respecto de la disponibilidad futura del recurso, especialmente para aquellos cultivos cuyo ciclo productivo se inicia en fechas cercanas a dicha actualización anual.

Debido a que los criterios utilizados por la Junta de Vigilancia no se encuentran formalizados mediante reglas explícitas de decisión, sino que corresponden a una evaluación conjunta de múltiples variables hidrológicas y meteorológicas, se optó por desarrollar un \textbf{modelo de regresión lineal múltiple} capaz de aproximar matemáticamente las relaciones observadas entre dichas variables y el porcentaje de desmarque registrado históricamente. 

La regresión lineal múltiple constituye una metodología ampliamente utilizada cuando se busca modelar una variable respuesta a partir de múltiples predictores, permitiendo estimar la contribución conjunta de cada variable explicativa y obtener una función analítica fácilmente integrable dentro del modelo de simulación (\citeref{james2021}{James et al., 2021}; \citeref{montgomery2021}{Montgomery et al., 2021}). Bajo este enfoque, el modelo no pretende reproducir el criterio subjetivo de los especialistas, sino aproximar empíricamente el comportamiento histórico de las decisiones adoptadas por la Junta de Vigilancia a partir de la información hidrológica disponible.

Para la construcción del modelo se utilizaron antecedentes hidrológicos y meteorológicos históricos provenientes principalmente de la red de estaciones CEAZAMet, la plataforma SWEET para el monitoreo de nieve y los boletines hidrológicos publicados por la Dirección General de Aguas (DGA), integrando variables representativas del estado de la cuenca durante cada temporada de análisis.

Previamente al proceso de pretratamiento de los datos, se recopiló y organizó la información hidrológica y meteorológica correspondiente a las distintas temporadas de estudio. La Tabla~\ref{tab:variables_modelo} presenta las variables consideradas en el proceso de modelación, junto con su descripción y unidad de medida.

\begin{figure}[H]
    \centering
    \vspace{0.3cm}
    [FIGURA PENDIENTE: Flujo metodológico seguido para construir el modelo de regresión (desde datos de entrada hasta modelo final)]
    \caption{Flujo metodológico de la regresión predictiva del desmarque}
    \label{fig:flujo_desmarque}
\end{figure}

\vspace{0.2cm} 

\begin{table}[H]
\centering
\caption{Nomenclatura y descripción de las variables utilizadas en el análisis.}
\label{tab:variables_modelo}
\begin{tabular}{lll}
\hline
\textbf{Variable} & \textbf{Descripción} & \textbf{Unidad} \\
\hline
Desmarque & Porcentaje de desmarque registrado para la temporada. & \% \\
RioElquiProm & Caudal promedio de la temporada en el río Elqui. & m$^{3}$/s \\
RioElquiMin & Caudal mínimo de la temporada en el río Elqui. & m$^{3}$/s \\
RioElquiMax & Caudal máximo de la temporada en el río Elqui. & m$^{3}$/s \\
PuclaroProm & Nivel promedio de la temporada del embalse Puclaro. & Mm$^{3}$ \\
LaLagunaProm & Nivel promedio de la temporada del embalse La Laguna. & Mm$^{3}$ \\
CoberturaProm & Cobertura nival promedio de la temporada. & km$^{2}$ \\
CoberturaMax & Cobertura nival máxima de la temporada. & km$^{2}$ \\
PrecipitacionTotal & Precipitación total acumulada de la temporada. & mm \\
TempProm & Temperatura media de la temporada  en la estación Vicuña. & $^\circ$C \\
\hline
\end{tabular}
\end{table}

Inicialmente, se intentó pronosticar el desmarque mediante estadísticas descriptivas de las principales variables hidrológicas, considerando sus valores mínimo, promedio y máximo, mientras que para las variables meteorológicas se utilizaron la precipitación acumulada y la temperatura media. Si bien esta aproximación permitía describir el comportamiento observado durante la temporada, generaba un conjunto amplio de variables derivadas de una misma serie temporal, aumentando la \textbf{redundancia} de la información y generando un importante \textbf{sobreajuste} en los primeros modelos (\citeref{kuhn2019}{Kuhn y Johnson, 2019}).

Posteriormente, durante el proceso de refinamiento del modelo y tras reuniones técnicas con especialistas del Departamento de Hidrología del Centro de Estudios Avanzados en Zonas Áridas (CEAZA), se identificó y revisó el informe de \citeref{rodhos2015}{RODHOS (2015)} como un antecedente técnico relevante para la cuenca. Cabe precisar que el modelo de regresión ya se encontraba en desarrollo a partir de los antecedentes recopilados previamente con la Junta de Vigilancia; por ello, la revisión de este informe no constituyó el origen de la metodología, sino que actuó como un antecedente técnico que enriqueció y fortaleció la selección de variables, motivando la incorporación explícita del volumen del Embalse La Laguna, predictor de cabecera que no había sido considerado inicialmente. Asimismo, las recomendaciones de los especialistas del CEAZA y los criterios de dicho informe técnico confirmaron la necesidad de representar cada temporada mediante descriptores sintéticos que condensaran el comportamiento hídrico sin incrementar la redundancia (\citeref{mckee1993}{McKee et al., 1993}; \citeref{vicente2010}{Vicente-Serrano et al., 2010}).

Este enfoque permite reducir la dimensionalidad del problema, disminuir la redundancia entre variables y obtener un conjunto de predictores más representativo para el ajuste del modelo de regresión. Asimismo, la construcción de descriptores sintéticos constituye una estrategia ampliamente utilizada en problemas de modelación hidrológica y aprendizaje automático cuando se busca preservar la información característica de una serie temporal sin incrementar innecesariamente la complejidad del modelo (\citeref{kuhn2019}{Kuhn \& Johnson, 2019}). Siguiendo esta recomendación, se desarrolló un indicador capaz de integrar el nivel medio y la variabilidad observada durante cada temporada, definido por:

\[
I=
\bar{x}
\left(
\frac{\bar{x}-x_{\min}}
{x_{\max}-x_{\min}}
\right),
\]

donde $\bar{x}$ corresponde al valor promedio de la temporada, mientras que $x_{\min}$ y $x_{\max}$ representan los valores mínimo y máximo observados, respectivamente. Cuando no existe variabilidad durante el período ($x_{\max}=x_{\min}$), el indicador adopta directamente el valor promedio.

Este procedimiento permitió representar cada temporada mediante un único descriptor por variable hidrológica, conservando la información más representativa del período, reduciendo la dimensionalidad del problema y obteniendo un conjunto de predictores más compacto y robusto para el ajuste del modelo de regresión.

La utilización de un único descriptor por temporada facilitó además la interpretación de los resultados y disminuyó la redundancia entre predictores, favoreciendo la estabilidad del modelo y una mejor capacidad de generalización frente a nuevas temporadas hidrológicas, objetivo ampliamente recomendado en la literatura de modelación predictiva (\citeref{kuhn2019}{Kuhn \& Johnson, 2019}).


Posteriormente, se evaluaron distintos modelos mediante procedimientos de selección de variables hacia adelante (\textit{Forward Selection}), hacia atrás (\textit{Backward Elimination}) y selección bidireccional (\textit{Stepwise}), seleccionándose finalmente el modelo con mejor capacidad predictiva:

\vspace{0.2cm}

\[
\hat{D}=1.2601
+0.0020\,Puclaro
-0.0033\,Cobertura
-0.0710\,Temperatura
+0.0013\,LaLaguna
\]
\begin{figure}[H]
    \centering
    \includegraphics[width=0.6\textwidth]{fig/Recta regresión.png}
    \caption{Recta modelo de regresión con selección hacia atrás}
    \label{fig:recta_regresion}
\end{figure}
\vspace{0.2cm}

El modelo obtenido alcanzó una \textbf{elevada capacidad predictiva}, con un coeficiente de determinación de $R^2=0.9886$ y un coeficiente de determinación ajustado de $R^2_{aj}=0.9773$. Asimismo, presentó un \textbf{error cuadrático medio} (RMSE) de $0.01026$ y un \textbf{sesgo prácticamente nulo} ($Bias=0.0000$), lo que indica que las predicciones no presentan sobreestimación ni subestimación sistemática.

\vspace{0.2cm}

Con el fin de verificar la \textbf{validez estadística} del modelo, se evaluaron los principales supuestos de la regresión lineal. Los resultados obtenidos mediante las pruebas de \textbf{Jarque--Bera} ($p=0.6202$), \textbf{Breusch--Pagan} ($p=0.5096$) y \textbf{Durbin--Watson} ($p=0.5264$) no evidenciaron incumplimientos significativos de los supuestos de normalidad, homocedasticidad e independencia de los residuos, respectivamente, permitiendo utilizar el modelo como base para la generación de escenarios.

Finalmente, a partir de esta estimación se construyen \textbf{cinco escenarios de disponibilidad hídrica} que representan la incertidumbre asociada al porcentaje de desmarque. En la implementación del modelo, el usuario define un \textbf{valor base de desmarque} $d_0$ correspondiente a su mejor estimación del porcentaje esperado —la cual puede apoyarse en la predicción $\hat{D}$ del modelo de regresión—, y los escenarios se generan variando dicho valor en un \textbf{paso fijo} $\Delta d$ en torno a ese valor central:

\vspace{0.2cm}

\[
d(i)=d_0+i\cdot \Delta d,
\qquad
i\in\{-2,-1,0,1,2\},
\qquad
\Delta d \approx RMSE,
\]

\vspace{0.2cm}

donde $d_0$ es el valor base de desmarque, $\Delta d$ el \textbf{salto entre escenarios} \textemdash fijado en el orden del error de predicción ($RMSE$) del modelo de regresión\textemdash\ e $i$ el índice de escenario. \textit{Propaga la incertidumbre del modelo predictivo hacia la oferta hídrica: el escenario central ($i=0$) refleja la estimación base del regante, mientras que los escenarios extremos ($i=\pm 2$) representan temporadas significativamente más secas o más húmedas.}

\vspace{0.2cm}

Cada uno de estos escenarios genera una trayectoria distinta del porcentaje de desmarque durante la temporada agrícola. Dichas trayectorias son incorporadas como parámetro de entrada al modelo de dinámica de sistemas, modificando diariamente la disponibilidad efectiva de agua superficial y permitiendo evaluar el comportamiento del sistema productivo bajo distintos niveles de incertidumbre hídrica.

\subsection{Modelo geoespacial del Canal Bellavista}

De manera análoga al uso de un modelo predictivo de regresión lineal múltiple para estimar el porcentaje de desmarque, el modelo de Dinámica de Sistemas requiere de un cálculo espacial de la ubicación del regante para representar de forma hidráulicamente realista la disponibilidad neta en el predio. Con el propósito de superar la simplificación agregada tradicional en la modelación de pérdidas por conducción, el modelo de oferta incorpora este cálculo geoespacial acoplado a la dinámica continua de flujos. Este enfoque permite reproducir el comportamiento espacial observado en el Canal Bellavista, donde la disponibilidad efectiva del recurso hídrico disminuye progresivamente hacia el tramo final del trazado, afectando directamente a los agricultores de la sección inferior, denominados localmente ``coleros del canal''.

El procedimiento para estructurar este cálculo geoespacial de pérdidas consistió en un flujo de trabajo en las siguientes etapas:
\begin{enumerate}
    \item \textbf{Adquisición y edición de la planimetría oficial:} Se obtuvo el plano catastral oficial del Canal Bellavista en formato CAD (DWG), suministrado por la asesoría de ingeniería de la Junta de Vigilancia. Debido a que el dibujo vectorial constaba de segmentos inconexos que dificultaban la continuidad topológica, se editó el archivo en AutoCAD para integrar y unir las líneas de las bifurcaciones, estructurándolas por sección de flujo para posibilitar mediciones de longitud uniformes.
    \item \textbf{Conversión de formato geoespacial:} Utilizando el software \textit{Global Mapper}, la información vectorial en formato DWG editada fue convertida y georreferenciada a formato KML, proyectando las coordenadas UTM y preservando la jerarquía de la red del canal principal y sus derivaciones secundarias.
    \item \textbf{Ajuste y modelamiento de la red:} A través de las plataformas \textit{Google Earth} y \textit{Google Maps}, se realizaron correcciones y ajustes visuales sobre imágenes satelitales reales para modelar de manera precisa la distribución de los canales secundarios y los puntos de entrega en las compuertas prediales.
    \item \textbf{Reconstrucción geométrica y cálculo de distancias:} El archivo KML depurado fue procesado en Python empleando librerías especializadas en análisis de datos geoespaciales (\texttt{geopandas}, \texttt{shapely} y \texttt{pyogrio}), reconstruyendo la red de flujo y localizando los predios. A partir de las coordenadas del regante, se determina la distancia hidráulica real desde la bocatoma (km 0,1) hasta su punto de extracción.
    \item \textbf{Segmentación de infraestructura:} El trazado recorrido por el recurso es desagregado en: kilómetros de canal principal revestido y sin revestir, longitud recorrida en las bifurcaciones o canales secundarios al interior del valle y, finalmente, la distancia de conducción desde el canal secundario hasta el predio del regante.
    \item \textbf{Parámetros de pérdidas empíricos:} Los coeficientes de pérdidas por conducción provienen de los aforos experimentales detallados en el informe técnico oficial \textit{``Mejoramiento Canales Bellavista, Villalón y Buzeta: Informe General Canal Bellavista''}, elaborado por ARCADIS Chile para la Comisión Nacional de Riego (CNR, Proyecto 4184-1000-GA-INF-007). Específicamente, para los tramos iniciales sin revestir comprendidos entre el kilómetro 0,0 y el kilómetro 8,1 del canal, el informe técnico de la consultora documenta la medición de 4 segmentos hidráulicos con tasas de pérdidas por infiltración y conducción que oscilan entre 0,2\%/km y 5,5\%/km, lo que sustenta la alta tasa de pérdidas en la sección no revestida de la red.
\end{enumerate}

\begin{figure}[H]
    \centering
    \includegraphics[width=0.85\textwidth]{fig/DWG.png}
    \caption{Planimetría catastral original del Canal Bellavista en formato CAD (DWG).}
    \caption*{\footnotesize Fuente: Asesoría de Ingeniería, Junta de Vigilancia del Río Elqui y sus Afluentes.}
    \label{fig:plano_dwg}
\end{figure}



Para representar formalmente esta restricción en la simulación de la oferta, el modelo de Dinámica de Sistemas incorpora la determinación de una cota máxima o límite físico de caudal disponible en el predio ($Q_{max}$), el cual actúa como un techo dinámico que restringe la disponibilidad de agua según la ubicación del regante en la red y las condiciones de operación de la temporada.

El impacto de este techo hidráulico varía significativamente dependiendo de los derechos de aprovechamiento de aguas (acciones de agua) del agricultor y del porcentaje de desmarque vigente en la temporada. Para un agricultor con pocas acciones de agua (por ejemplo, 10 acciones) y operando bajo un escenario de desmarque severo (por ejemplo, un 10\% o 20\%), su caudal bruto nominal es muy bajo. En este caso, el flujo demandado es menor que la capacidad de conducción física y las pérdidas del canal, por lo que el caudal que llega a la compuerta predial no alcanza a chocar con la cota máxima del tramo. En consecuencia, este límite no representa una restricción activa en su simulación diaria y el agricultor recibe la totalidad de su asignación nominal mermada solo por pérdidas marginales.

Por el contrario, la cota máxima cobra especial relevancia y se convierte en una restricción crítica en escenarios donde los regantes poseen un mayor volumen de derechos (por ejemplo, 20 acciones de agua en lugar de 10) o cuando la temporada permite un desmarque superior (por ejemplo, un 50\% o más). Bajo estas condiciones de mayor caudal bruto nominal disponible, el flujo supera el límite físico del tramo o la bifurcación secundaria (afectado por la longitud hidráulica recorrida y el tipo de revestimiento). En este caso, la cota máxima actúa como un techo efectivo, restringiendo el caudal que efectivamente ingresa al predio. Por ejemplo, un agricultor con derechos nominales equivalentes a 20 acciones puede verse limitado a recibir un caudal efectivo equivalente a solo 15 acciones debido a la restricción física de la red hidráulica.

Esta modelación representa con alta fidelidad la realidad de la distribución de aguas en el sector agrícola de Pan de Azúcar, identificando que la posición geográfica (ser ``colero del canal'') impone restricciones estructurales severas. Al integrar los coeficientes hidráulicos de pérdidas por conducción a lo largo de toda la red, se puede identificar la tendencia de los caudales, de sus respectivas pérdidas y un total de pérdida acumulada que alcanza hasta el 73,8\% respecto de la bocatoma de cabecera (\citeref{arcadis2015}{ARCADIS, 2015}). Por lo tanto, la determinación de esta cota máxima representa un factor fundamental en el Sistema de Soporte a la Decisión (DSS), ya que permite al agricultor evaluar de manera realista la viabilidad de sus estrategias de siembra y el comportamiento de sus cultivos frente a estas restricciones de infraestructura.

\subsection{Integración del modelo de oferta hídrica}

La integración de la oferta hídrica acopla los tres componentes anteriores. El modelo predictivo de regresión determina el escenario de desmarque $d(i)$ para la temporada. Este valor alimenta al modelo en Dinámica de Sistemas, el cual calcula el caudal bruto diario en función de las acciones del regante. Finalmente, este flujo es limitado por el modelo geoespacial mediante la cota máxima de conducción física del canal principal y bifurcaciones secundarias, entregando el caudal efectivo neto disponible diariamente en la compuerta del predio.

\section{Construcción del modelo de demanda agrícola}

El modelo de demanda agrícola simula la respuesta productiva del predio a nivel operativo ante el calendario de oferta hídrica entregado por el canal.

\subsection{Representación del sistema productivo}

La simulación del sistema productivo se desarrolla mediante un modelo de \textbf{Simulación de Eventos Discretos (SED)}, el cual representa la evolución del sistema agrícola a través de la ocurrencia de eventos asociados a las distintas labores realizadas durante el ciclo fenológico de los cultivos. De acuerdo con \citeref{banks2010}{Banks et al. (2010)}, este paradigma permite representar sistemas donde el estado evoluciona mediante eventos discretos, modelando explícitamente entidades, recursos, actividades y restricciones operacionales, siendo especialmente apropiado para la evaluación de distintos escenarios de operación.

La estrategia de representación espacial utilizada en esta investigación surge como una inspiración conceptual de los principios de discretización ampliamente utilizados en el análisis numérico. En particular, la idea proviene del proceso de aproximación introducido por las sumas de Riemann, donde un dominio continuo es dividido en un conjunto finito de particiones con el propósito de aproximar una magnitud continua, mejorando progresivamente la precisión de la aproximación a medida que aumenta el número de subdivisiones (\citeref{burden2011}{Burden \& Faires, 2011}). Aunque el presente modelo no implementa integración numérica, adopta esta misma filosofía de discretización para representar un predio agrícola continuo mediante un conjunto finito de unidades espaciales relativamente homogéneas.

Este principio de descomposición constituye además la base de numerosos métodos numéricos empleados en ingeniería. Un ejemplo ampliamente conocido corresponde al Método de los Elementos Finitos, donde un dominio continuo es discretizado en elementos sobre los cuales se resuelve localmente el problema, integrando posteriormente dichas soluciones para aproximar el comportamiento global del sistema (\citeref{bathe1996}{Bathe, 1996}; \citeref{zienkiewicz2013}{Zienkiewicz et al., 2013}). Si bien la presente investigación no implementa un Método de los Elementos Finitos, traslada esta filosofía de discretización al contexto agrícola, permitiendo representar explícitamente la heterogeneidad espacial del predio mediante unidades independientes de simulación.

Esta aproximación difiere de la mayor parte de la literatura revisada en simulación agrícola, donde la subdivisión del terreno suele asumirse únicamente como una decisión de modelación sin una fundamentación matemática explícita. En contraste, el presente trabajo propone una discretización espacial sustentada en principios provenientes del análisis numérico y de la ingeniería, adaptándolos al modelamiento agrícola mediante Simulación de Eventos Discretos.

Inicialmente, el usuario define la superficie total del predio, la cual es dividida en un número determinado de \textbf{particiones} que representan unidades independientes de cultivo. Cada partición constituye una unidad de simulación sobre la cual se modela de forma independiente el desarrollo fenológico del cultivo, el balance hídrico, la utilización de recursos y el desempeño económico. Finalmente, los resultados obtenidos en cada partición son integrados para representar el comportamiento global del predio.

Desde el punto de vista computacional, el número de particiones introduce un compromiso entre fidelidad espacial y eficiencia de cálculo. Un mayor nivel de discretización permite representar con mayor precisión la variabilidad espacial del sistema agrícola; sin embargo, también incrementa el número de combinaciones de cultivo que deben evaluarse y, por consiguiente, el tiempo computacional requerido para la simulación. Este compromiso entre precisión y costo computacional constituye un principio ampliamente reconocido tanto en simulación como en análisis numérico (\citeref{banks2010}{Banks et al., 2010}; \citeref{burden2011}{Burden \& Faires, 2011}). Considerando que en explotaciones hortícolas de pequeña y mediana escala la subdivisión física del terreno suele realizarse en un número reducido de sectores de manejo relativamente homogéneos, el modelo limita la discretización a un máximo de cuatro particiones por hectárea, obteniendo un equilibrio adecuado entre representatividad espacial y eficiencia computacional.

\subsection{Modelación mediante eventos discretos}

Posteriormente, se identifican los cultivos factibles para la fecha de inicio de la simulación mediante un \textbf{filtro de estacionalidad}, descartando aquellas especies cuyo calendario agrícola no permite su establecimiento en el período seleccionado. De esta forma, el espacio de búsqueda se reduce únicamente a alternativas técnicamente viables.

Una vez definidos los cultivos disponibles, el modelo genera \textbf{todas las combinaciones posibles} de asignación sobre las particiones del terreno, las cuales son evaluadas individualmente mediante simulación. Si existen $n$ cultivos factibles y el terreno se divide en $k$ particiones, el número de alternativas generadas está dado por

\[
\binom{n+k-1}{k}.
\]

Para cada combinación se consideran los parámetros agronómicos específicos de cada cultivo, incluyendo la duración de las etapas fenológicas, profundidad radicular, coeficientes de cultivo duales ($K_{cb}$ y $K_e$), rendimiento esperado y costos de producción por hectárea obtenidos a partir de las fichas técnicas de ODEPA, actualizados mediante el Índice de Precios al Consumidor (IPC). Asimismo, el usuario define las propiedades hidráulicas del suelo, tales como la capacidad de campo ($CC$), el punto de marchitez permanente ($PMP$), el agua fácilmente aprovechable ($AFA$) y el agua total aprovechable ($ATA$), parámetros utilizados para representar el almacenamiento de agua en el perfil del suelo.

\begin{figure}[H]
    \centering
    \vspace{0.3cm}
    [FIGURA PENDIENTE: Diagrama de la Red de Petri que representa el flujo operacional del sistema agrícola]
    \caption{Red de Petri del flujo operacional del sistema agrícola predial}
    \label{fig:petri_demand}
\end{figure}

\subsection{Cálculo de requerimientos hídricos}

Durante cada día de simulación se calcula el \textbf{balance hídrico} del cultivo utilizando el procedimiento descrito en el Manual FAO~56, estimando la evapotranspiración del cultivo mediante el enfoque de \textbf{coeficientes duales}:

\[
ET_c=(K_{cb}+K_e)\cdot ET_0,
\]

donde $ET_0$ corresponde a la evapotranspiración de referencia obtenida a partir de la estación meteorológica del CEAZA y $K_{cb}$ y $K_e$ representan el coeficiente basal del cultivo y el coeficiente de evaporación del suelo, respectivamente.

La aplicación del riego no depende únicamente de la demanda del cultivo, sino también de la disponibilidad diaria de agua proporcionada por el modelo de dinámica de sistemas. En consecuencia, el modelo integra la programación de las llegadas de agua superficial, la disponibilidad de agua subterránea, la frecuencia de riego definida por el agricultor, las políticas de gestión del recurso y la capacidad de almacenamiento del estanque cuando éste se encuentra disponible.

Esta interacción permite representar de forma conjunta la oferta y la demanda hídrica, determinando diariamente el estado del cultivo, el estrés hídrico acumulado y la evolución del proceso productivo hasta la cosecha.
\subsection{Cálculo de indicadores productivos y económicos}

Una vez finalizada la simulación de cada alternativa productiva, se calculan los indicadores técnicos, productivos y económicos que permiten evaluar integralmente el desempeño de cada combinación de cultivos. En administración agrícola, la toma de decisiones no depende exclusivamente del rendimiento físico obtenido, sino de la interacción entre productividad, utilización de recursos, costos de producción, calidad comercial y rentabilidad esperada (\citeref{kay2019}{Kay et al., 2016}). Asimismo, investigaciones recientes muestran que incrementos en el rendimiento no necesariamente se traducen en mayores beneficios económicos, ya que el desempeño del sistema agrícola depende simultáneamente de factores productivos, comerciales y económicos (\citeref{michler2019}{Michler et al., 2018}; \citeref{kuznietsova2024}{Kuznietsova et al., 2024}; \citeref{bezat2024}{Bezat-Jarzębowska et al., 2024}). En consecuencia, el modelo calcula un conjunto de indicadores que sintetizan el comportamiento técnico y económico del sistema, constituyendo la base para comparar escenarios y seleccionar la alternativa productiva más conveniente.

\vspace{0.2cm}

En primer lugar, la \textbf{producción esperada} de cada partición se obtiene a partir del rendimiento simulado y del área cultivada:

\vspace{0.2cm}

\begin{equation}
Q = Y \cdot A,
\end{equation}

\vspace{0.2cm}

donde $Q$ es la producción [kg], $Y$ el rendimiento simulado [kg/ha] y $A$ el área cultivada [ha]. La producción física constituye el primer indicador obtenido a partir de la simulación, ya que representa el resultado directo del comportamiento fenológico del cultivo bajo las condiciones hídricas simuladas. Sin embargo, de acuerdo con la literatura de administración agrícola, este indicador por sí solo resulta insuficiente para evaluar el desempeño económico de una explotación, siendo necesario complementarlo con variables asociadas a calidad comercial, costos e ingresos (\citeref{kay2019}{Kay et al., 2016}).

\vspace{0.2cm}

Según el estrés hídrico acumulado, la producción se reparte entre las categorías comerciales de \textbf{primera calidad} ($P_1$), \textbf{segunda calidad} ($P_2$) y \textbf{pérdida} ($P_p$), que cumplen la identidad de conservación:

\vspace{0.2cm}

\begin{equation}
P_1 + P_2 + P_p = 100\%.
\end{equation}

\vspace{0.2cm}

La clasificación de la producción en distintas categorías comerciales responde a una práctica ampliamente utilizada en la comercialización hortícola, donde los estándares de calidad determinan el precio efectivamente recibido por el productor. En consecuencia, el rendimiento físico obtenido no representa necesariamente el volumen comercializable bajo las mejores condiciones de mercado, por lo que resulta necesario distinguir entre producción de primera calidad, segunda calidad y pérdidas (\citeref{kay2019}{Kay et al., 2016}).

\vspace{0.2cm}

Cada categoría capta una fracción distinta del valor comercial ($100\%$ la primera, $60\%$ la segunda y $0\%$ la pérdida), lo que define el \textbf{factor de captación del valor}:

\vspace{0.2cm}

\begin{equation}
FCV = 1.0\,P_1 + 0.60\,P_2 + 0.0\,P_p = P_1 + 0.60\,P_2,
\end{equation}

\vspace{0.2cm}

donde $P_1$ y $P_2$ son las proporciones de primera y segunda calidad. Con el propósito de justificar técnicamente la ponderación del $60\%$ asignada a la producción de segunda calidad ($0.60\,P_2$), se realizó un análisis empírico a partir de la base de datos histórica de la Oficina de Estudios y Políticas Agrarias (ODEPA). Mediante la comparación de las series de precios mensuales de las categorías comerciales de primera y segunda calidad disponibles para los cultivos hortícolas del estudio, se calculó el porcentaje promedio del valor de mercado capturado por el producto de segunda calidad respecto al de primera calidad, obteniéndose en promedio una relación de captación del $60\%$.

Conceptualmente, este indicador representa la fracción del valor potencial que logra capturar el productor una vez considerada la calidad comercial obtenida durante la simulación. Su formulación sigue la lógica utilizada en administración agrícola, donde el desempeño económico depende no sólo del volumen producido, sino también de la calidad alcanzada y del precio efectivamente recibido en el mercado (\citeref{kay2019}{Kay et al., 2016}). De esta forma, el FCV transforma la producción física simulada en producción económicamente valorizada, penalizando la segunda calidad y anulando completamente las pérdidas.

\vspace{0.2cm}

Con ello, el \textbf{ingreso real} de cada cultivo se calcula como:

\vspace{0.2cm}

\begin{equation}
Ingreso = Q \cdot Precio \cdot FCV,
\end{equation}

\vspace{0.2cm}

donde $Precio$ es el precio esperado del mes de cosecha [\$/kg]. Este cálculo integra simultáneamente el volumen producido, el precio esperado y la calidad comercial alcanzada, permitiendo representar de forma más realista el valor económico efectivamente obtenido por el agricultor. Este enfoque resulta consistente con investigaciones recientes que demuestran que el rendimiento físico no constituye por sí solo un buen indicador del desempeño económico, siendo necesario incorporar simultáneamente precios, calidad y condiciones de mercado para evaluar adecuadamente las alternativas productivas (\citeref{michler2019}{Michler et al., 2018}; \citeref{kuznietsova2024}{Kuznietsova et al., 2024}).

\vspace{0.2cm}

Finalmente, el desempeño económico se obtiene mediante el \textbf{ingreso neto}:

\vspace{0.2cm}

\begin{equation}
Ingreso_{neto} = Ingreso - Costo,
\end{equation}

\vspace{0.2cm}

donde $Costo$ es el costo total de producción [\$], actualizado mediante IPC. El ingreso neto constituye el principal indicador económico utilizado para comparar las distintas alternativas simuladas, al integrar simultáneamente los efectos de la producción, la calidad comercial, los precios de venta y los costos de producción. Este enfoque coincide con modelos recientes de simulación económica agrícola, los cuales utilizan la rentabilidad como criterio principal para evaluar escenarios y apoyar la toma de decisiones (\citeref{bezat2024}{Bezat-Jarzębowska et al., 2024}). Asimismo, diversos estudios muestran que maximizar el rendimiento no implica necesariamente maximizar el beneficio económico, debido a la interacción existente entre productividad, precios y costos de producción (\citeref{kuznietsova2024}{Kuznietsova et al., 2024}).

\vspace{0.2cm}

Adicionalmente, el modelo reporta \textbf{indicadores técnicos} del comportamiento de la simulación: la duración del ciclo productivo, el consumo total de agua, el estrés hídrico acumulado y la utilización del presupuesto,

\vspace{0.2cm}

\begin{equation}
T_{ciclo} = L_{ini}+L_{des}+L_{med}+L_{fin},
\qquad
V_{agua} = \sum_{t=1}^{T_{ciclo}} R_t,
\end{equation}

\begin{equation}
S_h = \sum_{t=1}^{T_{ciclo}} \big(1 - K_{s,t}\big),
\qquad
U_{ppto} = \frac{Costo}{B}\times 100,
\end{equation}

\vspace{0.2cm}

donde $L_{ini}$, $L_{des}$, $L_{med}$ y $L_{fin}$ corresponden a la duración de las etapas fenológicas [días], $R_t$ representa la lámina de riego aplicada en el día $t$, $K_{s,t}$ el coeficiente diario de estrés hídrico y $B$ el presupuesto disponible [\$]. Estos indicadores permiten caracterizar el comportamiento interno del sistema simulado desde una perspectiva técnica, facilitando la interpretación de las diferencias observadas entre escenarios. De acuerdo con \citeref{banks2010}{Banks et al. (2010)}, los indicadores de desempeño constituyen uno de los principales productos de un modelo de simulación, ya que permiten comparar políticas de operación, analizar el comportamiento del sistema e identificar los factores responsables de las diferencias entre alternativas. En conjunto, estos indicadores describen la extensión temporal del cultivo, el agua total aplicada, el estrés hídrico acumulado y la fracción del presupuesto utilizada, siendo consistentes con la formulación agronómica propuesta en el \citeref{allen1998}{Manual FAO~56} para la representación de los procesos de evapotranspiración y respuesta del cultivo al déficit hídrico.
\item \textbf{Modelación económica de los cultivos}

Con el propósito de realizar una evaluación económica consistente entre las distintas alternativas agrícolas, se utilizaron series históricas de precios mayoristas obtenidas desde la Oficina de Estudios y Políticas Agrarias (ODEPA). Como se observa en la Figura \ref{fig:precios_hortalizas}, los precios presentan un comportamiento estacional, por lo que la utilización de un único precio de referencia podría introducir sesgos dependiendo de la fecha de cosecha de cada cultivo.

\begin{figure}[h]
    \centering
    \includegraphics[width=0.85\textwidth]{fig/precios_hortalizas.png}
    \caption{Evolución histórica de precios de hortalizas y evidencia de estacionalidad.}
    \caption*{\footnotesize Fuente: Elaboración propia con información de ODEPA.}
    \label{fig:precios_hortalizas}
\end{figure}

Para eliminar el efecto de la inflación y hacer comparables los valores registrados en distintos años, todos los precios históricos fueron actualizados mediante el Índice de Precios al Consumidor (IPC), utilizando la misma base temporal empleada para los costos de producción.

Posteriormente, para cada cultivo se calculó un precio representativo mensual, obtenido como el promedio de todos los registros históricos correspondientes a un mismo mes (enero, febrero, marzo, ..., diciembre). De esta forma, el modelo dispone de un precio esperado para cada cultivo dependiendo del mes en que finaliza la cosecha.

Una vez concluida la simulación, el ingreso bruto de cada cultivo se obtiene utilizando la producción estimada y el precio promedio correspondiente al mes de cosecha. Posteriormente, dicho ingreso es ajustado mediante el Factor de Captación del Valor ($FCV$), el cual depende de la distribución de la producción entre primera calidad, segunda calidad y pérdidas.

Finalmente, el ingreso neto de cada alternativa se calcula descontando los costos de producción previamente actualizados mediante IPC, permitiendo comparar de manera consistente la rentabilidad de las distintas combinaciones de cultivo bajo un mismo poder adquisitivo y considerando explícitamente la estacionalidad de los precios agrícolas.

\item \textbf{Selección de alternativas y sistema de apoyo a la decisión}

Una vez evaluadas todas las combinaciones de cultivo, el modelo conserva únicamente las alternativas \textbf{factibles}, es decir, aquellas cuyo costo total no supera el presupuesto disponible $B$:

\vspace{0.2cm}

\begin{equation}
Costo_c \le B.
\end{equation}

\vspace{0.2cm}

\textit{Descarta las combinaciones económicamente inviables antes de la evaluación por desempeño.}

\vspace{0.2cm}

Con el propósito de incorporar la calidad de la producción al proceso de decisión, se define el \textbf{indicador de calidad compuesta}:

\vspace{0.2cm}

\begin{equation}
IC = 0.60\,P_{1} + 0.40\,P_{2},
\end{equation}

\vspace{0.2cm}

donde $P_1$ y $P_2$ son las proporciones de primera y segunda calidad [\%]. \textit{Resume en un solo valor la calidad comercial de la alternativa, ponderando más la primera calidad que la segunda.}

\vspace{0.2cm}

El desempeño global de cada alternativa se sintetiza en el indicador \textbf{Score}, que combina el margen económico con la calidad obtenida:

\vspace{0.2cm}

\begin{equation}
Score_c = Ingreso_{neto,c}\cdot \frac{IC_c}{100},
\qquad
Ingreso_{neto,c} = Ingreso_c - Costo_c.
\end{equation}

\vspace{0.2cm}

\textit{Prioriza alternativas que logran simultáneamente un elevado ingreso neto y una alta proporción de primera calidad; ante ingresos similares, la calidad comercial define el orden.}

\vspace{0.2cm}

Finalmente, la alternativa recomendada $c^{*}$ es aquella que \textbf{maximiza el Score} dentro del conjunto factible:

\vspace{0.2cm}

\begin{equation}
c^{*} = \operatorname*{arg\,max}_{c\,:\,Costo_c \le B} \; Score_c.
\end{equation}

\vspace{0.2cm}

\textit{Constituye el criterio de decisión del modelo, resuelto por enumeración exhaustiva sobre las $\binom{n+k-1}{k}$ combinaciones factibles.} Los resultados se presentan mediante un \textbf{sistema de apoyo a la decisión} que permite comparar escenarios de disponibilidad hídrica y analizar el comportamiento técnico, productivo y económico de cada combinación simulada.

\section{Formulación del balance hídrico}

La etapa de \textbf{simulación del proceso productivo} descrita en la sección anterior se sustenta en un \textbf{balance hídrico diario}, que constituye el motor de cálculo del modelo. A continuación se detalla su formulación, basada en el \textbf{Manual FAO~56}.

\vspace{0.2cm}

En lugar de seguir el contenido de agua de forma directa, el modelo rastrea el \textbf{agotamiento de la zona radicular} $D_r$ (déficit de agua respecto de la capacidad de campo) y, de forma acoplada, el \textbf{agotamiento de la capa superficial de evaporación} $D_e$.

\vspace{0.2cm}

La capacidad de almacenamiento del suelo queda definida por el \textbf{agua total aprovechable} ($ATA$) y el \textbf{agua fácilmente aprovechable} ($AFA$):

\vspace{0.2cm}

\begin{equation}
ATA = 1000\,(CC - PMP)\,Z_e,
\qquad
AFA = p\cdot ATA,
\end{equation}

\vspace{0.2cm}

donde $Z_e$ es la profundidad radicular y $p$ la fracción de agotamiento que el cultivo tolera sin sufrir estrés.

\vspace{0.2cm}

La \textbf{evapotranspiración del cultivo} se calcula con el enfoque de \textbf{coeficientes duales}, separando la transpiración de la evaporación del suelo:

\vspace{0.2cm}

\begin{equation}
ET_c(t) = K_s\,K_{cb}(t)\,ET_0(t)\,f(\theta) \;+\; E_s,
\qquad
E_s = K_e\,ET_0(t),
\end{equation}

\vspace{0.2cm}

donde $K_{cb}$ es el coeficiente basal del cultivo, $E_s$ la evaporación del suelo y $K_e$ el coeficiente de evaporación, acotado por $K_e=\min\!\big(K_r(K_{c\,max}-K_{cb}),\,f_{ew}K_{c\,max}\big)$.

\vspace{0.2cm}

El \textbf{coeficiente de estrés hídrico} $K_s$ reduce la transpiración cuando el agotamiento supera el umbral $AFA$:

\vspace{0.2cm}

\begin{equation}
K_s =
\begin{cases}
1, & D_r \le AFA,\\[4pt]
\dfrac{ATA - D_r}{ATA - AFA}, & AFA < D_r \le ATA.
\end{cases}
\end{equation}

\vspace{0.2cm}

Adicionalmente, se incorpora un \textbf{factor de retención no lineal} $f(\theta)=\left(1 - D_r/ATA\right)^{\alpha}$, que ajusta la transpiración según la textura del suelo mediante el exponente $\alpha$.

\vspace{0.2cm}

El estado del suelo se actualiza de forma \textbf{recursiva a escala diaria}, sumando la demanda evapotranspirativa y descontando los aportes de precipitación $P_t$ y riego $R_t$, con el agotamiento acotado a los límites físicos del suelo:

\vspace{0.2cm}

\begin{equation}
D_{r,t} = D_{r,t-1} + ET_c(t) - P_t - R_t,
\qquad 0 \le D_{r,t} \le ATA.
\end{equation}

\vspace{0.2cm}

De este agotamiento se deriva la \textbf{humedad relativa del suelo} $\theta_t = \left(1 - D_{r,t}/ATA\right)\times 100\%$, cuyos extremos corresponden al \textbf{punto de marchitez permanente} ($PMP$, cuando $D_r=ATA$) y a la \textbf{capacidad de campo} ($CC$, cuando $D_r=0$).

\vspace{0.2cm}

El riego $R_t$ no depende únicamente de la demanda del cultivo, sino de la disponibilidad diaria de agua entregada por el modelo de oferta hídrica, aplicándose según una \textbf{jerarquía de fuentes}: primero el \textbf{canal} (en día de turno), luego el \textbf{estanque} de acumulación y, finalmente, el \textbf{agua subterránea}. De esta forma, el modelo integra la oferta y la demanda hídrica día a día, determinando el estrés acumulado y la evolución del cultivo hasta la cosecha, y permitiendo comparar distintas \textbf{políticas de riego} bajo condiciones de incertidumbre hídrica.


\section{Integración multiparadigma}

La integración de los distintos paradigmas de simulación se basa en un mecanismo de paso de mensajes y sincronización de variables a nivel diario. El modelo de oferta de Dinámica de Sistemas calcula diariamente el flujo de agua disponible y lo comunica a la simulación de eventos discretos. Esta última actualiza el estado del balance hídrico del cultivo en cada una de las particiones del terreno y determina si se gatilla un evento de riego o de estrés, modificando la variable de acumulación del estanque o los consumos de agua subterránea.

\begin{figure}[H]
    \centering
    \vspace{0.3cm}
    [FIGURA PENDIENTE: Diagrama de integración conceptual entre la Dinámica de Sistemas y la Simulación de Eventos Discretos]
    \caption{Integración conceptual multiparadigma de oferta y demanda}
    \label{fig:integracion_conceptual}
\end{figure}

\section{Diseño experimental}

El diseño experimental contempla la simulación sistemática de combinaciones de cultivos comunes en la zona (tomates, lechugas, alcayotas y cebollines) distribuidos en particiones homogéneas del predio. Se evalúan cinco escenarios hidrológicos para el desmarque de agua ($RMSE = \pm 10\%$, $\pm 20\%$) cruzados con tres políticas operacionales de riego del agricultor: riego por turno tradicional, riego con estanque de acumulación nocturna y riego de apoyo mediante pozo de agua subterránea. Los indicadores resultantes permiten comparar la robustez económica y la eficiencia del uso del agua de cada alternativa agrícola.

\chapter{Diseño de implementación y desarrollo del software}

\section{Arquitectura del DSS}

La implementación del modelo se concibe como una \textbf{arquitectura modular} organizada en componentes encadenados, que reproducen el flujo de la simulación multiparadigma desde la adquisición de datos hasta la recomendación final. Cada módulo cumple una función específica y entrega su salida como insumo directo del siguiente, de manera que la recomendación producida por el sistema es trazable hasta los datos de entrada. Estos módulos pueden agruparse, según su naturaleza, en tres bloques funcionales: los de \textbf{entrada y decisión}, los \textbf{modelos analíticos y de simulación} y los \textbf{modelos de proceso y evaluación}. La Figura~\ref{fig:diseno_implementacion} sintetiza esta arquitectura.

\section{Módulo de entrada y predicción (regresión y KML)}

El primer bloque consolida los \textbf{datos de entrada} (información hidrológica, meteorológica, parámetros agronómicos de los cultivos, precios históricos ajustados por IPC y parámetros físicos del predio). A partir de la información histórica, el \textbf{modelo predictivo del desmarque} ejecuta la regresión lineal múltiple con mínimos cuadrados ordinarios para proyectar el desmarque $\hat{D}$. Adicionalmente, el procesamiento geoespacial del archivo KML reconstruye la red física y los tramos revestidos y sin revestir del Canal Bellavista para alimentar el cálculo espacial de pérdidas:
$$\text{CAD (DWG)} \to \text{Global Mapper} \to \text{KML} \to \text{Python (GeoPandas)} \to \text{Geometría hidráulica y distancias}$$

\section{Módulo de simulación (PySD y SimPy)}

El núcleo computacional del sistema ejecuta el acoplamiento de los paradigmas de simulación:
\begin{itemize}
    \item \textbf{Dinámica de sistemas (PySD):} Simula la oferta hídrica diaria del canal y el acuífero, integrando las restricciones de la cota máxima hidráulica ($Q_{max}$) obtenida del modelo geoespacial.
    \item \textbf{Eventos discretos (SimPy):} Simula diariamente el ciclo fenológico de los cultivos y la demanda acumulada en cada partición del terreno, controlando el balance de humedad en el suelo mediante el procedimiento de coeficientes duales de la FAO-56.
\end{itemize}

\section{Interfaz de usuario y visualización}

El \textbf{modelo económico} valoriza la producción física simulada aplicando precios mensuales y el factor de captación del valor según la calidad de la cosecha. Posteriormente, el componente DSS consolida los indicadores técnicos y económicos y los presenta en una \textbf{interfaz de usuario interactiva} desarrollada en Dash y Plotly. Esta interfaz gráfica permite al agricultor ingresar parámetros de entrada (como superficie, presupuesto y acciones de agua) y visualizar mediante gráficos interactivos y tablas ordenadas el desempeño proyectado para cada alternativa de siembra.

\begin{figure}[H]
    \centering
    \vspace{0.3cm}
    [FIGURA PENDIENTE: Pantallazo de la pantalla principal del DSS interactivo en Dash]
    \caption{Pantalla principal de la interfaz de usuario del DSS}
    \label{fig:interfaz_dss_main}
\end{figure}

\begin{figure}[H]
    \centering
    \vspace{0.3cm}
    [FIGURA PENDIENTE: Pantallazo del panel de configuración de parámetros de la simulación (suelo, presupuesto, acciones)]
    \caption{Panel de configuración de la simulación predial}
    \label{fig:interfaz_dss_config}
\end{figure}

\begin{figure}[H]
    \centering
    \vspace{0.3cm}
    [FIGURA PENDIENTE: Pantallazo del módulo de visualización interactiva de resultados y desempeño técnico-económico de cultivos]
    \caption{Visualización de resultados y comparación de alternativas de siembra}
    \label{fig:interfaz_dss_results}
\end{figure}

\section{Despliegue de la solución}

La aplicación DSS se diseñó para ser desplegada en un servidor web local o en la nube, garantizando su portabilidad y acceso multiplataforma mediante navegadores de internet. La arquitectura de código abierto de Python facilita este despliegue sin costos asociados de licencias de software comercial.

\begin{figure}[H]
    \centering
    \vspace{0.3cm}
    [FIGURA PENDIENTE: Flujo de procesamiento geoespacial desde plano CAD DWG, conversión KML hasta el cálculo de pérdidas e integración en Python]
    \caption{Flujo de procesamiento para el cálculo geoespacial de pérdidas en el Canal Bellavista}
    \label{fig:procesamiento_geoespacial_canal}
\end{figure}


\begin{figure}[h]
    \centering
    \includegraphics[width=\textwidth]{fig/Diagrama preliminar.png}
    \caption{Modelo de implementación del sistema}
    \caption*{\footnotesize Fuente: Elaboración propia.}
    \label{fig:diseno_implementacion}
\end{figure}

\section{Estrategia de validación}

La validación del modelo multiparadigma contempla dos dimensiones complementarias: una cuantitativa para el balance de humedad en suelo, y una cualitativa orientada a la evaluación de la lógica conceptual por parte de expertos y usuarios.

\subsection{Validación de balance de humedad y rendimiento}

La validación experimental tiene como objetivo evaluar la capacidad del modelo para representar la dinámica de la humedad del suelo bajo condiciones reales de operación. Para ello, se compararán las estimaciones generadas por el modelo de simulación con mediciones obtenidas directamente en un predio agrícola representativo, utilizando las mismas condiciones de suelo, clima y manejo del riego.

Como primera etapa, se realizará una caracterización física del suelo mediante ensayos de laboratorio, con el propósito de determinar los parámetros requeridos por el modelo de balance hídrico. Entre los análisis considerados se encuentran la determinación de la textura del suelo, densidad aparente, capacidad de campo, punto de marchitez permanente y demás propiedades hidráulicas necesarias para representar adecuadamente el comportamiento del suelo dentro del proceso de simulación.

Posteriormente, se instalará un sensor de humedad de suelo que permitirá medir el contenido volumétrico de agua durante el período de evaluación. La primera medición obtenida será utilizada como condición inicial del modelo, mientras que las mediciones posteriores permitirán registrar la evolución temporal de la humedad del suelo bajo condiciones reales.

\begin{figure}[H]
    \centering
    \includegraphics[width=0.40\textwidth]{fig/sensor_humedad.png}
    \caption{Sensor de humedad de suelo utilizado durante la validación experimental.}
    \label{fig:sensor_humedad}
\end{figure}

En forma paralela, el modelo de simulación será ejecutado utilizando los parámetros obtenidos en la caracterización del suelo, las variables meteorológicas registradas por la estación meteorológica del CEAZA y las ecuaciones de balance hídrico implementadas en el modelo. Como resultado, se obtendrá una estimación diaria de la humedad del suelo para el mismo período considerado en las mediciones experimentales.

La capacidad predictiva del modelo será evaluada comparando los valores simulados ($H_i^{sim}$) con los valores observados por el sensor ($H_i^{obs}$). Para cuantificar las diferencias entre ambas series se calcularán indicadores de error de simulación.

Finalmente, los indicadores obtenidos permitirán cuantificar el grado de concordancia entre el modelo y las mediciones experimentales. En caso de identificarse diferencias sistemáticas entre los valores simulados y observados, se evaluará la necesidad de calibrar los parámetros del balance hídrico. Este proceso se repetirá en predios agrícolas a modo de réplica con el fin de tener una validación cuantitativa robusta.

\subsection{Validación cualitativa y face validity}

Con el propósito de complementar la validación cuantitativa, se realizará una validación cualitativa basada en el concepto de \textit{face validity}, la cual consiste en evaluar la coherencia conceptual y operativa del modelo mediante la revisión de especialistas y actores vinculados al sistema agrícola. Esta validación considera la participación de profesionales del Centro de Estudios Avanzados en Zonas Áridas (CEAZA), investigadores del proyecto PTI Hortícola y agricultores de la zona de estudio, quienes analizarán si la representación de los procesos implementados en el modelo es consistente con el comportamiento observado en un sistema agrícola real.

\section{Extracción de Datos}


En esta sección se describe el proceso de recopilación, integración y preparación de los datos utilizados para la construcción del modelo de simulación. La investigación combina múltiples fuentes de información provenientes de organismos públicos y plataformas especializadas. En particular, se emplearon registros climáticos e hidrológicos obtenidos desde CEAZAMet y la Dirección General de Aguas (DGA), así como información económica proporcionada por la Oficina de Estudios y Políticas Agrarias (ODEPA), correspondiente a series históricas de precios y costos de producción agrícola. Adicionalmente, junto a ello se realizó una caracterización de los agricultores del sector de Pan de Azúcar mediante entrevistas, lo que permitió representar en el modelo perfiles de regantes deacuerdo a las condiciones reales de la zona de estudio.

\subsection{Datos climáticos provenientes de CEAZAMet}

Para caracterizar las condiciones climáticas del sector de Pan de Azúcar se utilizaron registros históricos de la estación meteorológica CEAZAMet ubicada en las coordenadas (-30.0746, -71.2389). Esta estación proporciona información de variables meteorológicas relevantes para la estimación de la demanda hídrica de los cultivos, entre ellas temperatura del aire, humedad relativa, evapotranspiración de referencia, precipitación, velocidad del viento y radiación solar.

Las series históricas fueron empleadas tanto en la construcción de la base de datos climática utilizada por el modelo de simulación como en la generación de variables predictoras para el modelo de estimación del nivel de desmarque del canal.

\subsection{Datos de cobertura nival}

Con el objetivo de mejorar la estimación del nivel de desmarque del canal, se incorporó información sobre la cobertura nival de la cuenca del río Elqui. Estos datos fueron obtenidos desde la plataforma SWEET-Coquimbo, la cual monitorea la evolución de la cobertura de nieve en distintos sectores de la región.

En este estudio se utilizaron los registros correspondientes al sector Algarrobal, considerando el periodo de estudio entre 2014 y la actualidad. La variable empleada corresponde al porcentaje de cobertura nival, utilizado como indicador de la disponibilidad futura del recurso hídrico.

\subsection{Datos hidrológicos de la Dirección General de Aguas}

Los datos hidrológicos fueron obtenidos a partir de los boletines hidrométricos publicados por la Dirección General de Aguas (DGA). En particular, se utilizaron las series históricas del nivel del embalse Puclaro y del caudal del río Elqui, variables incorporadas como predictores en el modelo de regresión utilizado para estimar el nivel de desmarque del canal.

\subsection{Datos económicos provenientes de ODEPA}

La información económica de los cultivos fue obtenida desde la Oficina de Estudios y Políticas Agrarias (ODEPA). Se utilizaron las series históricas de precios mayoristas de hortalizas y tubérculos correspondientes al período 2014 a 2025, expresadas en precios reales mediante la actualización de estos con el Índice de Precios al Consumidor (IPC).


Pero para poder generar una base se comparación robusta entre cultivos, se promediaron todos los registros correspondientes a un mismo mes del año ; es decir, todos los valores observados en enero, todos los de febrero y así sucesivamente hasta diciembre.

Este procedimiento permitió suavizar la variabilidad interanual de los precios y obtener una serie estacional representativa, proporcionando una base homogénea y comparable para la evaluación económica de las distintas alternativas de cultivo. 
 
 La Figura~\ref{fig:precio_promedio_mensual} presenta el comportamiento de los precios promedio mensuales obtenidos tras el ajuste por IPC. Dado que el modelo utiliza precios expresados en moneda de una fecha base, su aplicación en períodos futuros requerirá actualizar nuevamente las series históricas mediante el IPC correspondiente, con el fin de mantener la comparabilidad de los resultados económicos. 

\begin{figure}[H]
    \centering
    \includegraphics[width=\textwidth]{fig/precio_promedio_mensual.png}
    \caption{Precio promedio mensual de las hortalizas consideradas en el estudio, calculado a partir de las series históricas 2014--2025 ajustadas por IPC.}
    \label{fig:precio_promedio_mensual}
\end{figure}

Asimismo, se emplearon las fichas técnicas de costos de producción elaboradas por ODEPA, las cuales contienen información sobre rendimientos,es decir cuanto produce (unidades o kilos) por hectárea, requerimientos de insumos y costos asociados a distintos cultivos hortícolas. Esta información permitió estimar los costos de producción considerados en el modelo económico para cada partición de hectarea productiva.

\subsection{Datos del laboratorio geoquímico}


\subsection{Caracterización de los agricultores}

La caracterización de los agricultores se realizó mediante entrevistas con productores del sector de Pan de Azúcar. A partir de la información recopilada se identificaron las principales fuentes de abastecimiento de agua para riego y se definieron tres perfiles representativos de regantes:

\begin{itemize}
    \item \textbf{Perfil 1:} agricultores que disponen exclusivamente de acciones de agua del canal de riego.
    
    \item \textbf{Perfil 2:} agricultores que, además de acciones del canal, cuentan con un estanque de almacenamiento el cual les permite implementar sistemas de riego tecnificado.
    
    \item \textbf{Perfil 3:} agricultores que poseen acciones de aguas del canal, estanque de almacenamiento y ademas derechos de aprovechamiento de aguas subterráneas, lo que les proporciona una mayor flexibilidad para la gestión del recurso hídrico.
\end{itemize}

Adicionalmente, los agricultores fueron clasificados según su tamaño productivo utilizando los criterios establecidos en la legislación chilena. Para la categoría de pequeño productor se adoptó la definición establecida en el artículo 13 de la Ley Orgánica N$^\circ$18.910 del Instituto de Desarrollo Agropecuario (INDAP), que considera como pequeño agricultor a quien posee activos no superiores a 3.500 UF, explota una superficie de hasta 13 hectáreas de riego básico (HRB) y obtiene la mayor parte de sus ingresos de actividades agrícolas o silvoagropecuarias.

Por otra parte, la Ley N$^\circ$20.412 establece como productor mediano a aquel cuyas ventas anuales superan las 2.400 UF y no exceden las 25.000 UF. En consecuencia, los productores con ventas superiores a 25.000 UF fueron clasificados como grandes agricultores.

La caracterización obtenida permitió parametrizar el modelo de simulación incorporando diferencias en la disponibilidad de infraestructura de riego, fuentes de abastecimiento de agua y capacidad productiva de los distintos tipos de agricultores presentes en la zona de estudio.


\section{Funcionamiento de la Simulación Multiparadigma}

La simulación multiparadigma desarrollada constituye el núcleo computacional del sistema de apoyo a la toma de decisiones propuesto en esta investigación. Su objetivo es representar el comportamiento del sistema agrícola bajo distintos escenarios de disponibilidad hídrica, permitiendo analizar cómo las variaciones en la oferta de agua afectan la planificación de cultivos, la estrategia de riego y los resultados productivos y económicos del predio.

Para ello, el modelo integra dos paradigmas de simulación que representan fenómenos de distinta naturaleza. La Dinámica de Sistemas modela la evolución temporal de la oferta hídrica, considerando los procesos que determinan la disponibilidad de agua durante la temporada agrícola. Posteriormente, la Simulación de Eventos Discretos utiliza dicha disponibilidad para representar el funcionamiento del sistema productivo, modelando el desarrollo de los cultivos, las decisiones de riego y la utilización de las distintas fuentes de abastecimiento de agua.

Esta separación permite representar cada fenómeno mediante el paradigma más adecuado, manteniendo ambos modelos desacoplados desde el punto de vista computacional e integrados mediante el intercambio de información. Como resultado, el sistema es capaz de evaluar distintos escenarios de operación y generar indicadores técnicos, productivos y económicos que sirven como apoyo al proceso de planificación agrícola.

\subsection{Arquitectura General}

La arquitectura de la simulación multiparadigma se encuentra organizada en dos módulos principales que operan de forma secuencial. El primer módulo corresponde al modelo de oferta hídrica, implementado mediante Dinámica de Sistemas, cuya función es estimar la disponibilidad diaria de agua para distintos escenarios de operación. El segundo módulo corresponde al modelo de demanda agrícola, desarrollado mediante Simulación de Eventos Discretos, el cual utiliza la información generada por el primer módulo para simular el comportamiento de los cultivos durante toda la temporada agrícola.

Ambos módulos se encuentran desacoplados, comunicándose únicamente mediante un archivo intermedio denominado \texttt{CalendarioOferta.csv}. Este archivo almacena la oferta diaria de agua generada para cada escenario de simulación y constituye la entrada principal del modelo de demanda agrícola. Esta estrategia facilita la reutilización de ambos módulos de manera independiente, permitiendo modificar o recalibrar cualquiera de ellos sin alterar la lógica de funcionamiento del otro.

La Figura~\ref{fig:arquitectura_multiparadigma} presenta la arquitectura general del modelo desarrollado. En ella se observa el flujo de información desde las fuentes de datos y parámetros de entrada hacia el modelo de oferta hídrica, la generación del calendario de disponibilidad de agua y, posteriormente, su utilización por el modelo de demanda agrícola para la evaluación de escenarios y la generación de indicadores de desempeño.

% Insertar Figura de Arquitectura General

\subsection{Modelo de Oferta Hídrica Dinámica de Sistemas}

El primer componente de la simulación multiparadigma corresponde al modelo de oferta hídrica, desarrollado utilizando el paradigma de Dinámica de Sistemas. Este paradigma fue seleccionado debido a que la disponibilidad de agua constituye un proceso continuo cuya evolución depende de múltiples factores que interactúan a lo largo del tiempo, tales como la distribución de derechos de agua, las pérdidas durante la conducción, la programación de turnos de riego y las interrupciones por mantenimiento de la infraestructura hidráulica.

El objetivo de este módulo es estimar la cantidad de agua superficial disponible diariamente para cada agricultor durante el horizonte de planificación. Para ello, el modelo simula la evolución temporal del recurso hídrico considerando la operación del canal de riego y las condiciones definidas para cada escenario de simulación. La representación continua del sistema permite capturar cómo pequeñas variaciones en la disponibilidad de agua se acumulan y repercuten en los días posteriores de la temporada agrícola.

Para la representación matemática del caudal efectivo disponible diariamente en el predio del regante ($Q_{ef}$, en L/s), el modelo integra la formulación de derechos de agua nominales y desmarque, acoplada espacialmente con la eficiencia hidráulica de conducción del canal:

\begin{equation}
Q_{ef}(t) = \min\left( Q_{bruto}(t),\; Q_{max} \right)
\end{equation}

Donde el caudal bruto nominal $Q_{bruto}(t)$ (L/s) asociado a las acciones de agua se calcula en función de la dotación y de la tasa de ajuste temporal o desmarque ($d(t) \in [0,1]$) determinada para cada escenario de disponibilidad hídrica:

\begin{equation}
Q_{bruto}(t) = N_{acciones} \times V_{accion} \times d(t)
\end{equation}

El caudal máximo físicamente disponible ($Q_{max}$, en L/s) representa el límite superior impuesto por las pérdidas de conducción en la red. Este factor se estima a partir de la distancia recorrida sobre el canal principal Bellavista y la bifurcación secundaria utilizando la red geométrica procesada por el modelo geoespacial:

\begin{equation}
Q_{max} = Q_{inicial} \times \eta(L)
\end{equation}

donde $Q_{inicial}$ es el caudal aforado en la cabecera (1145 L/s) y $\eta(L)$ representa la eficiencia hidráulica total del recorrido de longitud $L$, la cual depende tanto de la distancia recorrida sobre el trazado principal y secundario como del tipo de infraestructura y revestimiento existente (ARCADIS, 2015):

\begin{equation}
\eta(L) = \eta_{canal}(d_{km}) \times \eta_{bif}(b_{km})
\end{equation}

La eficiencia acumulada sobre el canal principal Bellavista de $d_{km}$ de longitud se compone multiplicativamente tramo a tramo, donde cada tramo $i$ presenta pérdidas por conducción fijadas ($P_i$, \%) de acuerdo con el levantamiento de ingeniería de ARCADIS (2015):

\begin{equation}
\eta_{canal}(d_{km}) = \prod_{i=1}^{M} \left( 1 - \frac{P_i}{100} \times \phi_i(d_{km}) \right)
\end{equation}

donde $\phi_i(d_{km})$ representa la proporción del tramo $i$ recorrida por el agua:

\begin{equation}
\phi_i(d_{km}) = \begin{cases} 
0.0 & d_{km} \le Inicio_i \\
\frac{d_{km} - Inicio_i}{Fin_i - Inicio_i} & Inicio_i < d_{km} < Fin_i \\
1.0 & d_{km} \ge Fin_i
\end{cases}
\end{equation}

Una vez que el recurso es derivado a compuerta de una bifurcación o canal interior secundario de longitud $b_{km}$, experimenta una pérdida lineal con tasa de fricción $\beta$ (por defecto $0,5\%/\text{km}$):

\begin{equation}
\eta_{bif}(b_{km}) = \max\left(0.0,\; 1 - \beta \times b_{km}\right)
\end{equation}

A partir de este caudal efectivo predial, se calcula el volumen diario que alimenta la dinámica del stock de agua de riego disponible ($m^3$) durante cada turno, aplicando adicionalmente el factor de eficiencia por tecnificación del riego ($\eta_{tec}$), que representa la fracción del agua entregada en compuerta que efectivamente llega a la zona radicular según el método de aplicación utilizado:

\begin{equation}
\text{Oferta Neta }(m^3) = \frac{Q_{ef}(t) \times \text{Horas Turno} \times 3600}{1000} \times \eta_{tec}
\end{equation}

El parámetro $\eta_{tec}$ adopta valores distintos según el sistema de riego del agricultor: riego por goteo ($\eta_{tec}=0{,}95$), riego por aspersión ($\eta_{tec}=0{,}80$) y riego tradicional por superficie ($\eta_{tec}=0{,}60$). Este parámetro se almacena por regante en el archivo \texttt{metodo\_riego.csv} y se registra mediante la columna \texttt{metodo\_riego} en \texttt{regantes.csv}, permitiéndole a la interfaz de configuración actualizar automáticamente el factor correspondiente al seleccionar el tipo de riego.

Con el propósito de incorporar la incertidumbre propia de la gestión hídrica, el modelo genera múltiples escenarios de disponibilidad de agua a partir de diferentes valores de desmarque. Cada escenario representa una posible condición de operación para la temporada agrícola, permitiendo posteriormente evaluar el efecto de una mayor o menor disponibilidad de agua sobre el desempeño de los cultivos.

Como resultado de la simulación, el módulo genera un calendario diario de oferta hídrica para cada escenario evaluado. Esta información se almacena en el archivo \texttt{CalendarioOferta.csv}, el cual constituye la salida del modelo de Dinámica de Sistemas y la principal entrada para el módulo de Simulación de Eventos Discretos, encargado de representar el comportamiento de los cultivos.

\begin{figure}[H]
    \centering
    \vspace{0.3cm}
    [FIGURA PENDIENTE: Diagrama de stock and flow (acumulación y flujos) del modelo de Dinámica de Sistemas de la oferta hídrica predial]
    \caption{Diagrama de stock and flow del modelo de oferta hídrica}
    \label{fig:oferta_stock_flow}
\end{figure}

\subsection{Modelo de Demanda Agrícola mediante Simulación de Eventos Discretos}

El segundo componente de la arquitectura corresponde al modelo de demanda agrícola, implementado mediante Simulación de Eventos Discretos. Su función es representar el comportamiento operativo del sistema productivo utilizando la información generada por el modelo de oferta hídrica junto con un conjunto de archivos de configuración que caracterizan el entorno agrícola, de acuerdo con la lógica operacional de la Red de Petri del modelo predial (ver Figura~\ref{fig:petri_demand}). A diferencia del módulo de Dinámica de Sistemas, este componente no genera disponibilidad de agua, sino que utiliza la oferta calculada para simular el desarrollo de los cultivos y evaluar diferentes alternativas de planificación.

La principal entrada del modelo corresponde al archivo \texttt{CalendarioOferta.csv}, generado por el módulo de Dinámica de Sistemas. Este archivo almacena, para cada escenario y para cada día del horizonte de simulación, el volumen de agua superficial disponible en el canal, información que posteriormente es consultada durante la ejecución del modelo para determinar la disponibilidad efectiva de agua para riego.

Además de este archivo de integración, el modelo utiliza diversos archivos CSV que permiten parametrizar completamente la simulación. El archivo \texttt{regantes.csv} almacena la información de cada agricultor o unidad productiva, incluyendo la superficie del predio, número de particiones, frecuencia de riego, capacidad y nivel inicial del estanque, disponibilidad de extracción subterránea, coordenadas geográficas del predio (latitud y longitud) y el método de riego utilizado (columna \texttt{metodo\_riego}). Este último campo determina automáticamente el factor de eficiencia $\eta_{tec}$ a partir del archivo de referencia \texttt{metodo\_riego.csv}, el cual tabula los tres tipos disponibles: goteo (0{,}95), aspersión (0{,}80) y tradicional (0{,}60).

Las características de los cultivos se encuentran almacenadas en el archivo \texttt{data\_cultivos.csv}, el cual contiene la información necesaria para representar el desarrollo fenológico de cada especie, incluyendo la duración de sus etapas, coeficientes de cultivo y parámetros utilizados durante el balance hídrico. Complementariamente, el archivo \texttt{productividad\_cultivos.csv} incorpora la información económica de cada cultivo, como costos de producción, rendimiento esperado, precios de venta e ingresos proyectados según el período de cosecha.

Las condiciones ambientales utilizadas durante la simulación se obtienen desde el archivo \texttt{datosclima.csv}, el cual contiene registros diarios de variables meteorológicas empleadas para representar la demanda hídrica de los cultivos. Asimismo, el archivo \texttt{calendario\_siembra.csv} define la estacionalidad agrícola mediante una matriz binaria, donde cada columna representa un mes del año y los valores indican si un determinado cultivo puede o no establecerse durante dicho período. Esta estructura permite que el modelo descarte automáticamente aquellas alternativas de cultivo que no cumplen con las restricciones de calendario.

Una vez cargados todos los archivos de entrada, el motor de Simulación de Eventos Discretos utiliza esta información para representar el comportamiento diario del sistema agrícola. Durante la simulación, el modelo consulta permanentemente la disponibilidad de agua proveniente de \texttt{CalendarioOferta.csv}, las características de los cultivos, las condiciones climáticas y la configuración de cada regante, actualizando el estado de cada predio y registrando los indicadores técnicos y económicos correspondientes.

Como resultado, el módulo genera distintos archivos de salida que contienen la evolución de la simulación, indicadores de consumo de agua, utilización de las diferentes fuentes de abastecimiento, rendimiento de los cultivos, calidad de producción, costos, ingresos y utilidad obtenida para cada escenario evaluado. Adicionalmente, se generan reportes consolidados en formato HTML que permiten visualizar y comparar el desempeño de las alternativas analizadas.


\subsection{Integración entre Paradigmas}

La simulación multiparadigma se estructura mediante una arquitectura modular, donde los modelos de Dinámica de Sistemas y Simulación de Eventos Discretos operan como componentes independientes. Ambos módulos ejecutan sus procesos de forma desacoplada y se comunican exclusivamente mediante archivos de intercambio de información, manteniendo una separación entre la generación de la oferta hídrica y la simulación de la demanda agrícola. Esta arquitectura facilita el mantenimiento del sistema, la incorporación de nuevos escenarios y la actualización de cualquiera de los módulos sin afectar el funcionamiento del otro.

El flujo de ejecución comienza con la carga de los archivos de configuración que parametrizan la simulación. Entre ellos se encuentran \texttt{regantes.csv}, encargado de caracterizar cada unidad productiva; \texttt{data\_cultivos.csv} y \texttt{productividad\_cultivos.csv}, que contienen la información agronómica y económica de los cultivos; \texttt{datosclima.csv}, con las variables meteorológicas utilizadas durante la simulación; y \texttt{calendario\_siembra.csv}, que define mediante una matriz binaria los períodos permitidos para el establecimiento de cada cultivo. Estos archivos constituyen la base de información utilizada por el modelo durante toda la ejecución de la simulación.

Una vez cargados los datos de entrada, el módulo de Dinámica de Sistemas ejecuta la simulación de la oferta hídrica para cada uno de los escenarios definidos. Como resultado, genera el archivo \texttt{CalendarioOferta.csv}, el cual almacena la disponibilidad diaria de agua superficial durante todo el horizonte de simulación. Este archivo constituye el mecanismo de integración entre ambos paradigmas y representa el único intercambio directo de información entre los módulos.

Posteriormente, el modelo de Simulación de Eventos Discretos utiliza el contenido de \texttt{CalendarioOferta.csv} junto con los restantes archivos de configuración para representar el comportamiento operativo del sistema agrícola. Durante la simulación, el motor consulta diariamente la disponibilidad de agua registrada para cada escenario y la integra con la información correspondiente a los cultivos, condiciones climáticas y características de los regantes, permitiendo actualizar el estado del sistema y calcular los indicadores técnicos y económicos asociados a cada alternativa evaluada.

La integración entre ambos paradigmas se realiza mediante un flujo unidireccional de información, donde el modelo de Dinámica de Sistemas genera la disponibilidad de agua que posteriormente es utilizada por el modelo de Simulación de Eventos Discretos. De esta manera, cada escenario de oferta hídrica es evaluado bajo las mismas condiciones de operación, garantizando la consistencia de las simulaciones y permitiendo comparar objetivamente las distintas alternativas de planificación agrícola.

Finalmente, el modelo genera distintos archivos de salida con los resultados de la simulación, incluyendo indicadores técnicos, productivos y económicos para cada escenario evaluado, además de reportes consolidados en formato HTML que permiten analizar y comparar el desempeño de las distintas configuraciones agrícolas. Esta estrategia de integración permite que ambos paradigmas conserven su independencia conceptual y computacional, facilitando la escalabilidad del sistema y la incorporación de nuevas funcionalidades o modificaciones futuras sin alterar la arquitectura general de la plataforma.

\begin{figure}[H]
    \centering
    \vspace{0.3cm}
    [FIGURA PENDIENTE: Esquema del flujo de información e interacción entre el modelo de oferta y el de demanda dentro del DSS]
    \caption{Flujo de información e integración entre los modelos de oferta y demanda}
    \label{fig:integracion_paradigmas_flujo}
\end{figure}


\section{Interfaz del Sistema de Soporte a la Decisión (DSS)}

El Sistema de Soporte a la Decisión se materializa en una interfaz web desarrollada en Dash que estructura el flujo de trabajo en tres etapas secuenciales: configuración de parámetros, ejecución de la simulación y análisis de resultados. El diseño está orientado a que el agricultor opere el sistema sin requerir conocimientos técnicos de modelamiento, traduciendo la complejidad del modelo multiparadigma en una experiencia de uso directa y progresiva.

\subsection{Configuración de parámetros}

El proceso comienza en la pestaña \textbf{Módulo 1 — Canal}, donde el usuario configura los parámetros asociados a los derechos de agua superficial: número de acciones, tasa por acción, duración del turno, horizonte temporal de la simulación, porcentaje de desmarque inicial y final con su respectivo salto entre escenarios, y el calendario de paradas programadas del canal (Figura~\ref{fig:interfaz_oferta}).

\begin{figure}[H]
    \centering
    \includegraphics[width=\textwidth]{fig/Reportes/interfaz_parametros_oferta.png}
    \caption{Interfaz de configuración del Módulo 1 — Oferta Hídrica}
    \label{fig:interfaz_oferta}
\end{figure}

La pestaña \textbf{Módulo 2 — Regante} consolida los parámetros prediales del agricultor seleccionado: superficie en hectáreas, fecha de inicio de siembra, número de particiones de terreno, presupuesto disponible, propiedades del suelo según el estándar FAO-56 (capacidad de campo, punto de marchitez permanente y coeficiente de retención no lineal), derechos y stock de agua subterránea, y características del estanque predial. A su vez, el panel derecho muestra la ubicación georreferenciada del predio sobre el trazado del canal, junto con el caudal máximo calculable según la posición geográfica del regante (Figura~\ref{fig:interfaz_regante}).

\begin{figure}[H]
    \centering
    \includegraphics[width=\textwidth]{fig/Reportes/interfaz_parametros_regante.png}
    \caption{Interfaz de configuración del Módulo 2 — Regante y parámetros prediales}
    \label{fig:interfaz_regante}
\end{figure}

Finalmente, la pestaña \textbf{Cultivos} permite gestionar el catálogo de especies disponibles para la planificación. Cada cultivo puede ser expandido para editar sus parámetros económicos —costo por hectárea, rendimiento, precio mensual según estacionalidad y meses de siembra permitidos— así como sus coeficientes agronómicos FAO-56: duración de las fases fenológicas ($L_{ini}$, $L_{des}$, $L_{med}$, $L_{fin}$), coeficiente de cultivo ($K_c$) y coeficiente basal ($K_{cb}$). Es posible agregar nuevos cultivos mediante el formulario inferior de la misma pestaña (Figuras~\ref{fig:interfaz_cultivos} y \ref{fig:interfaz_cultivos2}).

\begin{figure}[H]
    \centering
    \includegraphics[width=\textwidth]{fig/Reportes/interfaz_parametros_cultivo.png}
    \caption{Pestaña de configuración de cultivos — vista general del catálogo}
    \label{fig:interfaz_cultivos}
\end{figure}

\begin{figure}[H]
    \centering
    \includegraphics[width=\textwidth]{fig/Reportes/interfaz_parametros_cultivo2.png}
    \caption{Edición de parámetros FAO-56 y económicos de un cultivo desplegado}
    \label{fig:interfaz_cultivos2}
\end{figure}

Una vez configurados todos los parámetros, el usuario presiona el botón \textbf{Guardar Todo} para persistir los valores en los archivos de configuración y, a continuación, \textbf{Simular Todo} para desencadenar la ejecución secuencial de ambos módulos.

\subsection{Modelo de oferta hídrica}

Tras iniciar la simulación, el modelo de Dinámica de Sistemas calcula la oferta hídrica para los cinco escenarios de desmarque definidos. La interfaz de animación del modelo —accesible de forma paralela— permite visualizar en tiempo real la evolución diaria del sistema hídrico: variables dinámicas como oferta neta, caudal efectivo, pérdida por posición y recarga subterránea, junto con los indicadores acumulados del ciclo de turnos y el histograma de oferta superficial por turno (Figura~\ref{fig:sim_ds}). Esta etapa genera el archivo \texttt{CalendarioOferta.csv}, que alimenta directamente el modelo de eventos discretos.

\begin{figure}[H]
    \centering
    \includegraphics[width=\textwidth]{fig/Reportes/simulación dinamica de sistemas ux ui.png}
    \caption{Interfaz del modelo de Dinámica de Sistemas durante la ejecución de la simulación hídrica}
    \label{fig:sim_ds}
\end{figure}

Finalizada la ejecución del modelo de oferta hídrica, el sistema de eventos discretos evalúa automáticamente todas las combinaciones factibles de cultivos por partición —considerando el presupuesto disponible— bajo cada uno de los cinco escenarios. El conjunto de corridas resultante constituye el espacio de decisión que el DSS pone a disposición del agricultor para su análisis.

\subsection{Análisis de resultados mediante el DSS}

Los resultados se presentan en el reporte interactivo, organizado en seis pestañas que permiten analizar el desempeño del sistema de forma progresiva.

El \textbf{Dashboard} (Figura~\ref{fig:dss_dashboard}) ofrece una tabla comparativa consolidada de todos los escenarios, mostrando margen total, ingreso, costos, producción, agua aplicada, déficit hídrico, utilización del presupuesto y calidad de primera promedio. Las celdas destacadas identifican visualmente el óptimo de cada indicador, permitiendo al agricultor identificar de forma inmediata el escenario más favorable sin necesidad de revisar el detalle individual.

\begin{figure}[H]
    \centering
    \includegraphics[width=\textwidth]{fig/Reportes/reporte_cultivo_dashboard.png}
    \caption{Dashboard del DSS: comparación consolidada de escenarios hídricos}
    \label{fig:dss_dashboard}
\end{figure}

La pestaña \textbf{Cultivos} (Figura~\ref{fig:dss_cultivos}) presenta, para cada escenario y partición de terreno, el cultivo que maximiza el margen bajo la restricción de presupuesto. Esta vista permite al usuario contrastar cómo varía la recomendación de siembra a medida que la disponibilidad hídrica cambia entre escenarios.

\begin{figure}[H]
    \centering
    \includegraphics[width=\textwidth]{fig/Reportes/reporte_cultivo_recomedación_por_escenario.png}
    \caption{Cultivos seleccionados por escenario y partición de terreno}
    \label{fig:dss_cultivos}
\end{figure}

La pestaña \textbf{Presupuesto} (Figura~\ref{fig:dss_presupuesto}) muestra, mediante un gráfico de barras apiladas, la proporción del presupuesto utilizado y restante en cada escenario en relación con el límite definido por el agricultor. Esta visualización permite evaluar la eficiencia en la asignación de recursos según la oferta hídrica proyectada.

\begin{figure}[H]
    \centering
    \includegraphics[width=\textwidth]{fig/Reportes/reporte_cultivo_comparacion_de_presupuesto.png}
    \caption{Comparación de la utilización del presupuesto por escenario}
    \label{fig:dss_presupuesto}
\end{figure}

La pestaña \textbf{Calidad} (Figura~\ref{fig:dss_calidad}) expone la distribución porcentual de primera calidad, segunda calidad y pérdidas para la combinación óptima de cada escenario. Los resultados se presentan en un gráfico de barras apiladas que permite comparar el impacto que tiene la variabilidad hídrica sobre la calidad cosechada, cuantificando el riesgo productivo asociado a cada nivel de desmarque.

\begin{figure}[H]
    \centering
    \includegraphics[width=\textwidth]{fig/Reportes/reporte_cultivo_calidad_por_escenario.png}
    \caption{Distribución de calidad de cosecha por escenario hídrico}
    \label{fig:dss_calidad}
\end{figure}

La pestaña \textbf{Corridas experimentales} (Figura~\ref{fig:dss_corridas}) presenta la totalidad de las combinaciones evaluadas para el escenario seleccionado, ordenadas por un índice de desempeño compuesto por margen y calidad de primera. El usuario puede identificar las corridas factibles dentro del presupuesto y seleccionar cualquier fila para acceder al detalle completo de esa combinación, incluyendo la Carta Gantt de siembra, los indicadores de rentabilidad y la distribución de calidad por partición (Figura~\ref{fig:dss_corridas2}).

\begin{figure}[H]
    \centering
    \includegraphics[width=\textwidth]{fig/Reportes/reporte_cultivo_analisis _corridas_experimentales.png}
    \caption{Vista general de corridas experimentales evaluadas por escenario}
    \label{fig:dss_corridas}
\end{figure}

\begin{figure}[H]
    \centering
    \includegraphics[width=\textwidth]{fig/Reportes/reporte_cultivo_analisis _corridas_experimentales_2.png}
    \caption{Detalle de una corrida seleccionada: Carta Gantt, rentabilidad y calidad por partición}
    \label{fig:dss_corridas2}
\end{figure}

Un tercer nivel de detalle (Figura~\ref{fig:dss_corridas3}) muestra el balance hídrico diario de la corrida seleccionada. En la sección superior se informa el estado del estanque y del acuífero subterráneo al inicio y al término de la simulación, junto con el presupuesto utilizado respecto al disponible. A continuación se presentan dos gráficos: la distribución de la oferta superficial por turno —distinguiendo riego directo, almacenamiento en estanque y pérdida sin absorber— y la distribución del agua aplicada por cultivo a lo largo del período simulado. Finalmente, la sección \textit{Simulación diaria por partición} expone, para cada partición de terreno, la evolución de la humedad del suelo respecto a la capacidad de campo y el punto de marchitez permanente, la evapotranspiración real frente a la potencial, y el agua aplicada desglosada por fuente (canal, estanque y subterráneo).

\begin{figure}[H]
    \centering
    \includegraphics[width=\textwidth]{fig/Reportes/reporte_cultivo_analisis _corridas_experimentales_3.png}
    \caption{Balance hídrico diario y simulación por partición de la corrida seleccionada}
    \label{fig:dss_corridas3}
\end{figure}

Por último, la pestaña \textbf{Explorar combinación} (Figura~\ref{fig:dss_explorar}) permite al agricultor simular libremente cualquier asignación de cultivos por partición bajo el escenario de su elección. Al presionar \textit{Simular combinación}, el sistema calcula en tiempo real los indicadores técnicos y económicos de esa configuración específica, incluyendo Carta Gantt, margen por partición, distribución de calidad y déficit hídrico, sin necesidad de recorrer la tabla de corridas precomputadas.

\begin{figure}[H]
    \centering
    \includegraphics[width=\textwidth]{fig/Reportes/reporte_cultivo_explorar_cultivo.png}
    \caption{Explorador libre de combinaciones de cultivos por escenario y partición}
    \label{fig:dss_explorar}
\end{figure}


\section{Validación de la Simulación}

\subsection{Por Expertos}
\chapter{Resultados y discusión metodológica}

Este capítulo presenta el análisis cuantitativo y cualitativo de los resultados preliminares obtenidos mediante la integración del modelo de regresión, la simulación multiparadigma y el análisis de la cota máxima geoespacial.

\section{Validación del modelo de regresión}

El ajuste del modelo predictivo del desmarque permitió obtener una ecuación robusta fundamentada en las condiciones físicas de la cuenca (embalses Puclaro y La Laguna, cobertura nival y temperatura media). Los indicadores de error ($RMSE = 0.01026$) y bondad de ajuste ($R^2 = 0.9886$) aseguran la validez estadística del modelo, verificando que la regresión no presenta sesgos sistemáticos. Esto permite generar escenarios probabilísticos realistas que representan adecuadamente la variabilidad interanual de los derechos de aprovechamiento de aguas para el agricultor.

\section{Resultados de la simulación multiparadigma}

La ejecución del modelo multiparadigma sobre las particiones prediales permite evaluar el comportamiento operativo y económico del sistema. Los resultados preliminares se analizan bajo cuatro casos de estudio representativos del valle de Pan de Azúcar:
\begin{itemize}
    \item \textbf{Caso 1: Pequeño agricultor con riego por turnos tradicional.} Evalúa un predio de baja escala expuesto a las frecuencias de llegada del canal sin capacidad de almacenamiento.
    \item \textbf{Caso 2: Agricultor mediano con estanque de acumulación nocturna.} Simula el impacto amortiguador del estanque sobre la continuidad del riego.
    \item \textbf{Caso 3: Agricultor con pozo de agua subterránea de apoyo.} Analiza el beneficio de contar con una fuente de respaldo frente a cortes programados del canal.
    \item \textbf{Caso 4: Regante al término del canal (colero del canal).} Evalúa el efecto extremo de la posición geográfica y las pérdidas por fricción acumulativas sobre la viabilidad de los cultivos.
\end{itemize}

\section{Evaluación de la cota máxima geoespacial}

La incorporación del catastro planimétrico en formato KML y los coeficientes hidráulicos oficiales de ARCADIS (2015) confirmaron que la localización física impone una restricción severa sobre el flujo de agua predial. Para agricultores con derechos nominales elevados en períodos de alto desmarque, la cota máxima de conducción física actúa como un techo dinámico que reduce el caudal efectivo real que ingresa a la compuerta. Este comportamiento reproduce el desequilibrio histórico observado en terreno y dota a la recomendación agrícola de un realismo físico ausente en los modelos de simulación convencionales.

\section{Propuesta de valor y soporte a la decisión}

Para analizar el funcionamiento teórico de la simulación planteada y validar su propuesta de valor, se presentó el diseño conceptual y los resultados preliminares ante profesionales especializados y actores relevantes:
\begin{itemize}
    \item \textbf{CEAZA:} Profesionales de esta institución validaron el planteamiento hidrológico y edafológico del balance de humedad en el suelo (PMP y capacidad de campo).
    \item \textbf{PTI Hortícola:} El agrónomo Luis Felipe Algo corroboró que la formulación y lógica operativa del DSS son consistentes con las necesidades de toma de decisiones del mercado regional, sugiriendo futuras convergencias para afinar la precisión en márgenes reducidos de consumo.
    \item \textbf{Comunidad de regantes:} La presentación del prototipo del DSS en mesas de comunidad con agricultores locales evidenció una excelente recepción cualitativa, rescatando la visualización interactiva de escenarios como una herramienta clave para reducir la incertidumbre antes de iniciar la siembra.
\end{itemize}

\chapter{Planificación del proyecto}

Este capítulo detalla la planificación temporal, hitos y cronograma general bajo el cual se estructuró el desarrollo de la investigación.

\section{Fases del proyecto}

El proyecto se organizó en cuatro fases encadenadas:
\begin{itemize}
    \item \textbf{Fase 0: Formulación del Proyecto.} Formalización del project charter y formulación del problema (ver Tabla~\ref{tab:fase0}).
    \item \textbf{Fase 1: Levantamiento de Datos y Calibración.} Recopilación de datos meteorológicos, edáficos, hidrológicos y entrevistas en terreno (ver Tabla~\ref{tab:fase1}).
    \item \textbf{Fase 2: Desarrollo y Simulación Multiparadigma.} Modelamiento continuo, de eventos discretos y acoplamiento en Python (ver Tabla~\ref{tab:fase2}).
    \item \textbf{Fase 3: Evaluación de Políticas y Análisis.} Análisis de escenarios, comparación estadística y entrega del manuscrito final (ver Tabla~\ref{tab:fase3}).
\end{itemize}

\begin{table}[h]
\centering
\begin{tabularx}{\textwidth}{|X|c|c|c|}
\hline
\textbf{Proceso/Actividad} & \textbf{Progreso} & \textbf{Inicio} & \textbf{Fin} \\
\hline
Segmentación de datos y generación de escenarios climáticos & 100\% & 20/04/26 & 26/04/26 \\
\hline
Modelado de Dinámica de Sistemas para la oferta hídrica & 100\% & 26/04/26 & 07/05/26 \\
\hline
Modelado de Eventos Discretos para la planificación operativa & 100\% & 07/05/26 & 13/05/26 \\
\hline
Consolidación de resultados preliminares y preparación de la defensa & 100\% & 19/05/26 & 08/06/26 \\
\hline
\textbf{Hito: Inicio de redacción manuscrito entregable final}  &  & 27/04/26 &  \\
\hline
\textbf{Hito: Presentación de Avance} & & \textbf{08/06/26} & \\
\hline
\end{tabularx}
\caption{Cronograma de actividades - Fase 2: Desarrollo y Simulación Multiparadigma}
\label{tab:fase2}
\end{table}

\begin{table}[h]
\centering
\begin{tabularx}{\textwidth}{|X|c|c|c|}
\hline
\textbf{Proceso/Actividad} & \textbf{Progreso} & \textbf{Inicio} & \textbf{Fin} \\
\hline
Parametrización de políticas de riego (Persecución de demanda, Optimización) & 100\% & 09/06/26 & 14/06/26 \\
\hline
Ejecución de múltiples réplicas bajo escenarios de desmarque y cortes de canal & 100\% & 15/06/26 & 21/06/26 \\
\hline
Análisis estadístico de resultados & 100\% & 22/06/26 & 28/06/26 \\
\hline
\textbf{Hito: Entrega de borrador de resultados} & & \textbf{28/06/26} & \\
\hline
\end{tabularx}
\caption{Cronograma de actividades - Fase 3: Evaluación de Políticas y Análisis}
\label{tab:fase3}
\end{table}

\begin{thebibliography}{}

\bibitem{}\hypertarget{ref:allen1998}{Allen, R. G., Pereira, L. S., Raes, D., \& Smith, M. (1998). \textit{Crop evapotranspiration: Guidelines for computing crop water requirements} (FAO Irrigation and Drainage Paper No. 56). Food and Agriculture Organization.}
\vspace{0.3cm} 
\bibitem{}\hypertarget{ref:banks2010}{Banks, J., Carson, J. S., Nelson, B. L., \& Nicol, D. M. (2010). \textit{Discrete-event system simulation} (5th ed.). Pearson.} 

\vspace{0.3cm} 
\bibitem{}\hypertarget{ref:bezat2024}{Bezat-Jarzębowska, A., Rembisz, W., \& Jarzębowski, S. (2024). Maintaining agricultural production profitability—A simulation approach to wheat market dynamics. Agriculture, 14(11), 1910.   
https://doi.org/10.3390/agriculture14111910}
\vspace{0.3cm} 
\bibitem{}\hypertarget{ref:bathe1996}{Bathe, K.-J. (1996). Finite element procedures. Prentice Hall.}
\vspace{0.3cm} 
\bibitem{}\hypertarget{ref:boulahya2023}{Boulahya, I., Kuper, M., Le Page, C., \& Rinaudo, J. D. (2023). An agent-based model to analyze the effects of farmers’ decision-making on water use and salinization in irrigated agriculture. \textit{Agricultural Water Management, 276}, 108041. https://doi.org/10.1016/j.agwat.2022.108041}

\vspace{0.3cm} 
\bibitem{}\hypertarget{ref:burden2011}{Burden, R. L., \& Faires, J. D. (2011). Numerical analysis (9th ed.). Brooks/Cole, Cengage Learning.}

\vspace{0.3cm} 
\bibitem{}\hypertarget{ref:cnr2015}{Comisión Nacional de Riego. (2015). \textit{Mejoramiento canales Bellavista, Villalón y Buzeta: Informe ejecutivo canal Bellavista}. Ministerio de Agricultura, Gobierno de Chile.}

\vspace{0.3cm} 
\bibitem{}\hypertarget{ref:dga2019}{Dirección General de Aguas. (2019). \textit{Boletín hidrométrico N°489}. Ministerio de Obras Públicas, Gobierno de Chile.}
\vspace{0.3cm} 
\bibitem{}\hypertarget{ref:dga2026}{Dirección General de Aguas. (2026). \textit{Boletín hidrométrico N°574}. Ministerio de Obras Públicas, Gobierno de Chile.}

\vspace{0.3cm} 
\bibitem{}\hypertarget{ref:dingre2020}{Dingre, S. K., \& Gorantiwar, S. D. (2020). Determination of the water requirement and crop coefficient values of sugarcane by field water balance method in semiarid region. \textit{Agricultural Water Management, 232}, 106042. https://doi.org/10.1016/j.agwat.2020.106042}

\vspace{0.3cm} 
\bibitem{}\hypertarget{ref:elfartassi2025}{El Fartassi, I., Milne, A. E., Metcalfe, H., El Alami, R., Diarra, A., Alonso-Chavez, V., Zawadzka, J., Waine, T. W., \& Corstanje, R. (2025). An agent-based model of farmer decision making: Application to shared water resources in arid and semi-arid regions. Agricultural Water Management, 310, 109357.
https://doi.org/10.1016/j.agwat.2025.109357}

\vspace{0.3cm} 
\bibitem{}\hypertarget{ref:jia2026}{Jia, X., \& Zhang, R. (2026). Land–Water Allocation, Yield Stability, and Policy Trade-Offs Under Climate Change: A System Dynamics Analysis. Systems, 14(4), 412. 
https://doi.org/10.3390/systems14040412}

\vspace{0.3cm} 
\bibitem{}\hypertarget{ref:kuschel2020}{Kuschel-Otárola, M., Rivera, D., Holzapfel, E., Schütze, N., Neumann, P., \& Godoy-Faúndez, A. (2020). Simulation of water-use efficiency of crops under different irrigation strategies. \textit{Water, 12}(10), 2930. https://doi.org/10.3390/w12102930}


\vspace{0.3cm} 
\bibitem{}\hypertarget{ref:michler2019}{Michler, J. D., Tjernström, E., Verkaart, S., \& Mausch, K. (2019). Money matters: The role of yields and profits in agricultural technology adoption. American Journal of Agricultural Economics, 101(3), 710–731. https://doi.org/10.1093/ajae/aay050}

\vspace{0.3cm} 
\bibitem{}\hypertarget{ref:munyaka2024}{Munyaka, J.-C. B., Gallay, O., Hlal, M., Mutandwa, E., \& Chenal, J. (2024). Optimizing the sweet potato supply chain in Zimbabwe using discrete event simulation: A focus on production, distribution, and market dynamics. Sustainability, 16(21), 9166. https://doi.org/10.3390/su16219166}

\vspace{0.3cm} 
\bibitem{}\hypertarget{ref:lara2014}{Lara, A. (2014). \textit{Metodología de la investigación en simulación}. Pearson Educación.}

\vspace{0.3cm} 
\bibitem{}\hypertarget{ref:li2023}{Li, X., Zhang, J., Cai, X., Huo, Z., \& Zhang, C. (2023). Simulation-optimization based real-time irrigation scheduling: A human-machine interactive method enhanced by data assimilation. \textit{Agricultural Water Management, 276}, 108059. https://doi.org/10.1016/j.agwat.2022.108059}

\vspace{0.3cm} 
\bibitem{}\hypertarget{ref:codigo1981}{Ministerio de Justicia. (1981). \textit{Código de Aguas}. Biblioteca del Congreso Nacional de Chile.}

\vspace{0.3cm} 
\bibitem{}\hypertarget{ref:odepa}{ODEPA. (s.f.). \textit{Series de precios: Series de tiempo}. https://aplicativos.odepa.gob.cl/series-precios/series-tiempo}

\vspace{0.3cm} 
\bibitem{}\hypertarget{ref:olivera2025}{Olivera Ruiz, G., \& Pérez Espinoza, G. K. (2025). Sensores y aplicaciones de IoT en entornos agrícolas: Mapeo sistemático. \textit{Revista Invecom}.}

\vspace{0.3cm} 
\bibitem{}\hypertarget{ref:kay2019}{Kay, R. D., Edwards, W. M., \& Duffy, P. A. (2019). Farm management (9th ed.). McGraw-Hill Education.}

\vspace{0.3cm} 
\bibitem{}\hypertarget{ref:kuhn2019}{Kuhn, M., \& Johnson, K. (2019). Feature Engineering and Selection: A Practical Approach for Predictive Models. CRC Press.}

\vspace{0.3cm} 
\bibitem{}\hypertarget{ref:arcadis2015}{ARCADIS Chile. (2015). \textit{Mejoramiento Canales Bellavista, Villalón y Buzeta: Informe General Canal Bellavista} (Proyecto N° 4184-1000-GA-INF-007). Comisión Nacional de Riego (CNR), Ministerio de Agricultura, Gobierno de Chile.}
\bibitem{}\hypertarget{ref:kuznietsova2024}{Kuznietsova, T., Krasovska, Y., Lesniak, O., Harnaha, O., \& Podlevska, O. (2024). Case study on Ukrainian farm gross margin and direct cost analysis of wheat production. BIO Web of Conferences, 114, 01030. 
https://doi.org/10.1051/bioconf/202411401030}

\vspace{0.3cm} 
\bibitem{}\hypertarget{ref:mckee1993}{McKee, T. B., Doesken, N. J., \& Kleist, J. (1993). The relationship of drought frequency and duration to time scales.}
\vspace{0.3cm} 
\bibitem{}\hypertarget{ref:vicente2010}{Vicente-Serrano, S. M., Beguería, S., \& López-Moreno, J. I. (2010). A multiscalar drought index sensitive to global warming: The Standardized Precipitation Evapotranspiration Index. Journal of Climate.}
\vspace{0.3cm} 
\bibitem{}\hypertarget{ref:rivera2025}{Rivera, D., \& Donoso, G. (Eds.). (2025). \textit{Droughts in Chile: Impacts, monitoring, and adaptation policies}. Springer. https://doi.org/10.1007/978-3-031-85040-0}

\vspace{0.3cm} 
\bibitem{}\hypertarget{ref:rodhos2015}{RODHOS Asesorías y Proyectos Ltda. (2015). \textit{Desarrollo de modelos de gestión de recursos hídricos superficiales y subterráneos para las cuencas de las provincias de Elqui, Limarí y Choapa}. RODHOS Asesorías y Proyectos Ltda.}
\vspace{0.3cm} 
\bibitem{}\hypertarget{ref:un2023}{United Nations. (2023). \textit{The United Nations World Water Development Report 2023: Partnerships and cooperation for water}. UNESCO.}

\vspace{0.3cm} 
\bibitem{}\hypertarget{ref:james2021}{James, G., Witten, D., Hastie, T., \& Tibshirani, R. (2021). An introduction to statistical learning: With applications in R (2nd ed.). Springer. https://doi.org/10.1007/978-1-0716-1418-1}

\vspace{0.3cm} 
\bibitem{}\hypertarget{ref:montgomery2021}{Montgomery, D. C., Peck, E. A., \& Vining, G. G. (2021). Introduction to linear regression analysis (6th ed.). John Wiley \& Sons.}

\vspace{0.3cm} 
\bibitem{}\hypertarget{ref:montanari2013}{Montanari, A., Young, G., Savenije, H. H. G., Hughes, D., Wagener, T., Ren, L., Koutsoyiannis, D., Cudennec, C., Toth, E., Grimaldi, S., Blöschl, G., Sivapalan, M., Beven, K., Gupta, H., Hipsey, M., Schaefli, B., Arheimer, B., Boegh, E., Schymanski, S. J., ... Di Baldassarre, G. (2013). "Panta Rhei—Everything Flows": Change in hydrology and society—The IAHS Scientific Decade 2013–2022. Hydrological Sciences Journal, 58(6), 1256–1275. https://doi.org/10.1080/02626667.2013.809088}

\vspace{0.3cm} 
\bibitem{}\hypertarget{ref:kuhn2019b}{Kuhn, M., \& Johnson, K. (2019). Feature engineering and selection: A practical approach for predictive models. CRC Press.}

\vspace{0.3cm} 
\bibitem{}\hypertarget{ref:arcadis2015b}{ARCADIS Chile. (2015). \textit{Mejoramiento Canales Bellavista, Villalón y Buzeta: Informe General Canal Bellavista} (Proyecto N° 4184-1000-GA-INF-007). Comisión Nacional de Riego (CNR), Ministerio de Agricultura, Gobierno de Chile.}

\vspace{0.3cm}
\bibitem{}\hypertarget{ref:zienkiewicz2013}{Zienkiewicz, O. C., Taylor, R. L., \& Zhu, J. Z. (2013). \textit{The finite element method: its basis and fundamentals}. Elsevier.}

\vspace{0.3cm}
\bibitem{}\hypertarget{ref:jvre2025}{Junta de Vigilancia del Río Elqui. (2025, 12 de septiembre). \textit{JVRE aprueba por tercer año consecutivo un desmarque del 15\% para asegurar la temporada 2025–2026}. Junta de Vigilancia del Río Elqui y sus Afluentes. https://www.rioelqui.cl/jvre-aprueba-por-tercer-ano-consecutivo-un-desmarque-del-15-para-asegurar-la-temporada-2025-2026/}

\end{thebibliography}
\end{document}
\bibitem{}
 
McKee, T. B., Doesken, N. J., \& Kleist, J. (1993). The relationship of drought frequency and duration to time scales.
\vspace{0.3cm} 
\bibitem{}

Vicente-Serrano, S. M., Beguería, S., \& López-Moreno, J. I. (2010). A multiscalar drought index sensitive to global warming: The Standardized Precipitation Evapotranspiration Index. Journal of Climate.
\vspace{0.3cm} 
\bibitem{}

Rivera, D., \& Donoso, G. (Eds.). (2025). \textit{Droughts in Chile: Impacts, monitoring, and adaptation policies}. Springer. https://doi.org/10.1007/978-3-031-85040-0

\vspace{0.3cm} 

\bibitem{}
RODHOS Asesorías y Proyectos Ltda. (2015). \textit{Desarrollo de modelos de gestión de recursos hídricos superficiales y subterráneos para las cuencas de las provincias de Elqui, Limarí y Choapa}. RODHOS Asesorías y Proyectos Ltda.
\vspace{0.3cm} 
\bibitem{}
United Nations. (2023). \textit{The United Nations World Water Development Report 2023: Partnerships and cooperation for water}. UNESCO.

\vspace{0.3cm} 
\bibitem{}
James, G., Witten, D., Hastie, T., \& Tibshirani, R. (2021). An introduction to statistical learning: With applications in R (2nd ed.). Springer. https://doi.org/10.1007/978-1-0716-1418-1 


\vspace{0.3cm} 
\bibitem{}
Montgomery, D. C., Peck, E. A., \& Vining, G. G. (2021). Introduction to linear regression analysis (6th ed.). John Wiley \& Sons. 

\vspace{0.3cm} 
\bibitem{} 

Montanari, A., Young, G., Savenije, H. H. G., Hughes, D., Wagener, T., Ren, L., Koutsoyiannis, D., Cudennec, C., Toth, E., Grimaldi, S., Blöschl, G., Sivapalan, M., Beven, K., Gupta, H., Hipsey, M., Schaefli, B., Arheimer, B., Boegh, E., Schymanski, S. J., ... Di Baldassarre, G. (2013). "Panta Rhei—Everything Flows": Change in hydrology and society—The IAHS Scientific Decade 2013–2022. Hydrological Sciences Journal, 58(6), 1256–1275. https://doi.org/10.1080/02626667.2013.809088 


\vspace{0.3cm} 
\bibitem{} 

Kuhn, M., \& Johnson, K. (2019). Feature engineering and selection: A practical approach for predictive models. CRC Press.

\vspace{0.3cm} 
\bibitem{} 
ARCADIS Chile. (2015). \textit{Mejoramiento Canales Bellavista, Villalón y Buzeta: Informe General Canal Bellavista} (Proyecto N° 4184-1000-GA-INF-007). Comisión Nacional de Riego (CNR), Ministerio de Agricultura, Gobierno de Chile.


\end{thebibliography}
\end{document}

Kuhn, M., & Johnson, K. (2019). Feature Engineering and Selection: A Practical Approach for Predictive Models. CRC Press.

\vspace{0.3cm} 
\bibitem{}

Kuznietsova, T., Krasovska, Y., Lesniak, O., Harnaha, O., & Podlevska, O. (2024). Case study on Ukrainian farm gross margin and direct cost analysis of wheat production. BIO Web of Conferences, 114, 01030. 
https://doi.org/10.1051/bioconf/202411401030 

\vspace{0.3cm} 
\bibitem{}
 
McKee, T. B., Doesken, N. J., & Kleist, J. (1993). The relationship of drought frequency and duration to time scales.
\vspace{0.3cm} 
\bibitem{}

Vicente-Serrano, S. M., Beguería, S., & López-Moreno, J. I. (2010). A multiscalar drought index sensitive to global warming: The Standardized Precipitation Evapotranspiration Index. Journal of Climate.
\vspace{0.3cm} 
\bibitem{}

Rivera, D., \& Donoso, G. (Eds.). (2025). \textit{Droughts in Chile: Impacts, monitoring, and adaptation policies}. Springer. https://doi.org/10.1007/978-3-031-85040-0

\vspace{0.3cm} 

\bibitem{}
RODHOS Asesorías y Proyectos Ltda. (2015). \textit{Desarrollo de modelos de gestión de recursos hídricos superficiales y subterráneos para las cuencas de las provincias de Elqui, Limarí y Choapa}. RODHOS Asesorías y Proyectos Ltda.
\vspace{0.3cm} 
\bibitem{}
United Nations. (2023). \textit{The United Nations World Water Development Report 2023: Partnerships and cooperation for water}. UNESCO.

\vspace{0.3cm} 
\bibitem{}
James, G., Witten, D., Hastie, T., & Tibshirani, R. (2021). An introduction to statistical learning: With applications in R (2nd ed.). Springer. https://doi.org/10.1007/978-1-0716-1418-1 


\vspace{0.3cm} 
\bibitem{}
Montgomery, D. C., Peck, E. A., & Vining, G. G. (2021). Introduction to linear regression analysis (6th ed.). John Wiley & Sons. 

\vspace{0.3cm} 
\bibitem{} 

Montanari, A., Young, G., Savenije, H. H. G., Hughes, D., Wagener, T., Ren, L., Koutsoyiannis, D., Cudennec, C., Toth, E., Grimaldi, S., Blöschl, G., Sivapalan, M., Beven, K., Gupta, H., Hipsey, M., Schaefli, B., Arheimer, B., Boegh, E., Schymanski, S. J., ... Di Baldassarre, G. (2013). "Panta Rhei—Everything Flows": Change in hydrology and society—The IAHS Scientific Decade 2013–2022. Hydrological Sciences Journal, 58(6), 1256–1275. https://doi.org/10.1080/02626667.2013.809088 


\vspace{0.3cm} 
\bibitem{} 

Kuhn, M., & Johnson, K. (2019). Feature engineering and selection: A practical approach for predictive models. CRC Press.

\vspace{0.3cm} 
\bibitem{} 
ARCADIS Chile. (2015). \textit{Mejoramiento Canales Bellavista, Villalón y Buzeta: Informe General Canal Bellavista} (Proyecto N° 4184-1000-GA-INF-007). Comisión Nacional de Riego (CNR), Ministerio de Agricultura, Gobierno de Chile.
